% !TeX root = ../../Book2.tex
% 纲要标题：## 第15章　有限维代数、根滤过与投射覆盖
% 纲要位置：计划纲要.md，第 1525 行。

\chapter{有限维代数、根滤过与投射覆盖}
\label{b2:ch:15}
\BookChapterNumber{b2:ch:15}

\begin{chapterabstract}
有限维代数的 Jacobson 根是幂零理想，而模掉根的商是半单代数。两项结论使一个模可以按根的作用分层，但这些半单层之间仍可能有不分裂的扩张。本章用根与基座刻画 Loewy 列，计算三角矩阵代数的具体层次，再通过幂等元的相容提升构造投射覆盖。覆盖的存在性、唯一性和不可分解投射模的分类均在此证明。最后选取基本角代数，以半单商给出箭图的顶点，以 \(J/J^2\) 给出箭头。
\end{chapterabstract}

\begin{learninggoals}
\begin{itemize}
\item 区分降链稳定与根幂零，准确说明 Nakayama 引理在证明中的作用。
\item 计算模的根、基座和 Loewy 层，并辨认这些层之间是否分裂。
\item 完成幂零理想模的单个及正交幂等元组提升，利用多余核构造投射覆盖。
\item 说明不可分解投射模与简单模的对应，保留一般基域上自同态除环的作用。
\item 构造基本满角代数，按左模约定理解根第一层的箭头方向。
\end{itemize}
\end{learninggoals}

令 \(A\) 为域 \(k\) 上的二阶上三角矩阵代数。严格上三角矩阵组成理想 \(J\)，且 \(J^2=0\)；商 \(A/J\) 只看见两个对角位置，因而同构于 \(k\times k\)。商映射保留两个对角位置，却把非零的矩阵单位 \(E_{12}\) 送到零，也丢去了对角元与它相乘的作用。若研究一个 \(A\) 模，先模掉 \(J\) 能读出顶层，却无法独自恢复各层如何接合。

\ProseParagraphBreak

沿 \(J\) 的幂逐层观察，最先消失的乘法信息就能成为层之间关系的线索。若要从半单商中的简单对象回到原代数，还需找到足够小、又能映满目标的投射模。幂等元的提升和投射覆盖正是为这种回返服务；在三角矩阵例子中，两个对角幂等元与非对角位置已经把顶点和连接它们的箭头显露出来。

\ProseParagraphBreak

本章需要\cref{b2:thm:projective-characterizations}中的投射模刻画、\cref{b2:thm:artin-wedderburn}中的半单代数结构，以及\cref{b2:thm:noncommutative-nakayama}。角环及其与自同态环的关系使用\cref{b2:prop:corner-endomorphism}。除幂等元提升处明确讨论一般环外，\(A\) 始终是域 \(k\) 上的有限维结合含幺代数，模默认是幺性左模；对模的有限性要求在陈述中写明。

\ProseParagraphBreak
有限维代数及投射覆盖可参阅 Etingof 等人的《Introduction to Representation Theory》\cite[\S 3.5; \S\S 9.1--9.2]{EtingofEtAl2011RepresentationTheory}。该书第9章明确假设基域代数闭。本章的主要构造保留一般基域，简单模的自同态可能是大于 \(k\) 的除环；在计数 Hom 空间维数时尤其要区别这两种假设。根幂零性的主要证明采用第十四章的 Nakayama 引理，与该书 Proposition 3.5.3 的合成列证明可以相互参照。

% !TeX root = ../../Book2.tex
% 纲要标题：### 15.1 根的幂零性
% 纲要位置：计划纲要.md，第 1527 行。

\section{根的幂零性}
\label{b2:sec:1501}

有限维性使左理想的升降链都停止，但“链停止”还不等于“某个幂为零”。根滤过的各层可以通过半单商来理解，而要让滤过最终到达零，还需将有限维性与非交换 Nakayama 引理合起来。始终令 \(A\) 为域 \(k\) 上的有限维结合含幺代数，记 \(J=J(A)\)。

% 纲要要点（供目录核对）：
%
% * 有限维代数的 Jacobson 根
% * 半单商、根滤过与伴随分次代数
% * 幂零根和一般根的区别；区分逐元素幂零与统一混合乘积为零，说明伴随分次乘法保留的层次
% * 对有限维代数 \(A\) 的降链 \(J(A)^n\) 先取稳定项，再用第14.2节 Nakayama 引理推出其为零；不把 Artin 性本身当作根幂零性的全部证明
%

\subsection{有限维代数的 Jacobson 根}
\label{b2:subsec:1501:01}

回顾\cref{b2:def:jacobson-radical,b2:prop:radical-simple-annihilators}，\(J(A)\) 既是全部极大左理想的交，也是全部简单幺左模的共同零化子，因而是双边理想。特别地，对每个简单左模 \(S\)，都有 \(JS=0\)。后面各层被 \(J\) 零化的计算，使用的正是这个性质。

\ProseParagraphBreak
在 \(A\) 有限维的当前条件下，每个简单幺左 \(A\) 模都是 \(A\) 的商，因而有限维。每个有限生成幺左 \(A\) 模也有限维：有限个生成元给出满射 \(A^n\to M\)。反过来，若 \(m_1,\ldots,m_d\) 是 \(M\) 的 \(k\) 基，则 \(\sum_i c_im_i=\sum_i(c_i1_A)m_i\)，其中 \(c_i1_A\in A\)，所以它们也是 \(A\) 模生成元。这说明此时有限生成与有限 \(k\) 维等价，并不说明 \(A\) 模的最少生成元数等于 \(k\) 维数。

\begin{definition}\label{b2:def:nilpotent-ideal}
设 \(R\) 为含幺环、\(I\) 为其双边理想。对正整数 \(N\)，\(I^N\) 是全部 \(N\) 个 \(I\) 中元素的乘积所生成的加法子群。若某个正整数 \(N\) 使 \(I^N=0\)，则称 \(I\) 为\term[mi-ling-lixiang]{幂零理想}[nilpotent ideal]，满足该式的最小正整数称为其\term[miling-zhishu]{幂零指数}[nilpotency index]。
\end{definition}

第十四章出现的 nil ideal 只要求 \(\forall x\in I\,\exists n_x\ge1\)，使 \(x^{n_x}=0\)。幂零理想则要求 \(\exists N\ge1\,\forall x_1,\ldots,x_N\in I\)，都有 \(x_1\cdots x_N=0\)。后者除了使用一个共同指数，还要控制不同元素的混合乘积；仅考察同一个元素反复相乘，不能直接得到 \(I^N=0\)。

\begin{proposition}\label{b2:prop:nilpotent-ideal-in-radical}
任意含幺环的幂零双边理想都包含在其 Jacobson 根中。
\end{proposition}
\begin{proof}
若 \(I^N=0\)，则每个 \(x\in I\) 都满足 \(x^N=0\)，故 \(I\) 是 nil 双边理想。\cref{b2:cor:nil-ideal-radical}已证明这种理想包含于 Jacobson 根；它允许各元素使用不同的幂次，因而也适用于这里的共同幂次 \(N\)。
\end{proof}

例如 \(A=k[t]/(t^m)\)、\(m\ge1\)，其理想 \((\bar t)\) 幂零，且商为域 \(k\)。它因此恰是 \(J(A)\)：包含关系由上命题得到，而商的根为零，再用\cref{b2:prop:radical-quotient}即得反向包含。这里的根不是全部零因子的任意集合，而是由理想及简单模刻画的特定对象。

\subsection{半单商与根滤过}
\label{b2:subsec:1501:02}

\begin{proposition}\label{b2:prop:finite-algebra-semisimple-quotient}
设 \(k\) 为域、\(A\) 为有限维结合含幺 \(k\) 代数，\(J=J(A)\) 为其 Jacobson 根。商代数 \(A/J\) 是半单代数。每个被 \(J\) 零化的有限维幺左 \(A\) 模都是半单模。
\end{proposition}
\begin{proof}
由\cref{b2:prop:radical-quotient}，\(J(A/J)=0\)。商代数的左正则模有限维，故具有有限长度；其极大子模正是极大左理想，交为零。应用\cref{b2:prop:finite-length-radical-zero}，可知这个正则模半单，所以 \(A/J\) 是半单代数。

若 \(JM=0\)，定义 \((a+J)m=am\)。由于 \(J\) 在 \(M\) 上为零，此定义与代表元无关，使 \(M\) 成为幺左 \(A/J\) 模。它是某个有限自由 \((A/J)^n\) 的商。有限直和与商保持半单性，见\cref{b2:prop:semisimple-subquotients}，故 \(M\) 半单；其 \(A/J\) 子模与 \(A\) 子模是同一批子空间。
\end{proof}

由\cref{b2:thm:artin-wedderburn}，可以写
\begin{equation}\label{b2:eq:finite-algebra-semisimple-quotient}
A/J\cong\prod_{i=1}^s M_{n_i}(D_i),
\end{equation}
其中各 \(D_i\) 是有限维 \(k\) 除代数。暂不假设 \(k\) 代数闭，所以不能把 \(D_i\) 一律写成 \(k\)。

\ProseParagraphBreak
只要 \(J\) 在一个幺左模上作用为零，\(A\) 的作用就能且只能通过 \(A/J\) 给出；反过来，每个幺左 \(A/J\) 模沿商映射限制标量，都被 \(J\) 零化。两种作用给出的子模和模同态完全相同，所以简单性及半单性也相同。这里 \(A/J\) 半单，是根为零与有限长度共同带来的结论；对任意环，仅有 Jacobson 根为零不足以推出半单性。

\ProseParagraphBreak
\(J\) 的各个幂给出\term[gen-lvguo]{根滤过}[radical filtration]
\[
A\supseteq J\supseteq J^2\supseteq\cdots.
\]
每个 \(J^r/J^{r+1}\) 都被 \(J\) 零化，因为 \(J J^r=J^{r+1}\)，所以由\cref{b2:prop:finite-algebra-semisimple-quotient}，它作为左模半单。类似地，对任何有限维左模 \(M\)，商 \(J^rM/J^{r+1}M\) 也是半单模。各层半单并不说明原模是这些层的直和；层与层之间的扩张可以不分裂。

\ProseParagraphBreak
若只把这些商空间逐个列出，\(A\) 的乘法就不在数据中了。但 \(J^rJ^s\subseteq J^{r+s}\)，所以第 \(r\) 层和第 \(s\) 层的代表元相乘，仍能在第 \(r+s\) 层读出一个商类。伴随分次代数把这种乘法也记录下来。

\begin{definition}\label{b2:def:associated-radical-algebra}
设 \(k\) 为域、\(A\) 为有限维结合含幺 \(k\) 代数，\(J=J(A)\)，并约定 \(J^0=A\)。在分次向量空间
\[
\operatorname{gr}_J(A)=\bigoplus_{r\ge0}J^r/J^{r+1}
\]
上，对非负整数 \(r,s\) 及 \(a\in J^r\)、\(b\in J^s\) 定义
\[
(a+J^{r+1})(b+J^{s+1})=ab+J^{r+s+1}.
\]
所得分次 \(k\) 代数称为\term[gen-lvguo-bansui-fenci-daishu]{根滤过的伴随分次代数}[associated graded algebra of the radical filtration]。
\end{definition}
\symidx{\(\operatorname{gr}_J(A)\)：根滤过的伴随分次代数}

乘积不依赖代表元：若分别把 \(a,b\) 改动一个 \(J^{r+1},J^{s+1}\) 中的元素，新增各项都在 \(J^{r+s+1}\) 中。结合律由 \(A\) 的乘法继承，单位为 \(1+J\)，而次数按和相加。这个乘法只记录 \(ab\) 模 \(J^{r+s+1}\) 的值；若 \(ab\) 实际落在更深的 \(J^{r+s+1}\) 中，即使 \(ab\ne0\)，它在第 \(r+s\) 次的乘积也为零。因此较深层的修正项会被丢掉，不能仅凭伴随分次代数恢复 \(A\) 的全部乘法。

\ProseParagraphBreak
在 \(A=k[t]/(t^m)\) 中，\(J^r\) 的基为 \(\bar t^r,\ldots,\bar t^{m-1}\)，\(0\le r<m\)。这个例子恰有与滤过相容的分次结构：令 \(u\) 的次数为一，映射 \(k[u]/(u^m)\to\operatorname{gr}_J(A)\) 将 \(u\) 送到 \(\bar t+J^2\)，将 \(u^r\) 送到第 \(r\) 层的 \(\bar t^r+J^{r+1}\)。它保持乘法，并把各次的一维基对应起来，因而是分次 \(k\) 代数同构。这里能恢复截断多项式代数，依赖于这一具体乘法，不能作为一般情形的结论。每个非零层又是一维简单模 \(k\)，但当 \(m>1\) 时正则模不是这些层的直和：\(\bar t\) 在任何这样的简单模直和上均作用为零，而在正则模上并非如此。

\subsection{幂零根与一般根}
\label{b2:subsec:1501:03}

对有限维代数，本节将证明 \(J\) 是幂零理想。这个结论比“\(J\) 中每个元素各自幂零”更强，因为同一个整数要使任意若干根元素的乘积都为零。它也不意味着所有幂零元素都在根中。

\ProseParagraphBreak
例如 \(M_2(k)\) 是半单代数，故根为零，但矩阵单位 \(E_{12}\ne0\) 满足 \(E_{12}^2=0\)。幂零的单个元素生成的双边理想未必幂零；这里 \(E_{21}E_{12}=E_{22}\)，所以它生成的双边理想已经含有非零幂等元。由于 \(E_{22}^n=E_{22}\ne0\) 对每个正整数 \(n\) 都成立，这个双边理想的各个幂都非零。\cref{b2:prop:nilpotent-ideal-in-radical}的假设是整个双边理想幂零。

\ProseParagraphBreak
若去掉有限维性，Jacobson 根可以完全不幂零。例如形式幂级数环 \(k[\![t]\!]\) 中，常数项非零的级数可以逐次解出逆级数的系数，因而可逆；常数项为零的级数不可能可逆。故非单位元恰为理想 \((t)\)，它是唯一极大理想，等于 Jacobson 根。但 \(t^n\notin(t^{n+1})\)，所以 \((t)^n\supsetneq(t)^{n+1}\) 对所有 \(n\ge0\) 都成立，降链根本没有稳定项。Nakayama 引理仍适用；只是没有 \(J^{n+1}=J^n\) 这个等式，就不能把它用于稳定项来推出零。

\subsection{根幂零性的 Nakayama 证明}
\label{b2:subsec:1501:04}

\begin{theorem}\label{b2:thm:finite-algebra-radical-nilpotent}
设 \(k\) 为域、\(A\) 为有限维结合含幺 \(k\) 代数，\(J=J(A)\) 为其 Jacobson 根，则 \(J\) 幂零。若 \(A\ne0\)，可以取正整数 \(N\le\dim_k A\) 使 \(J^N=0\)。因此 \(J\) 是 \(A\) 中最大的幂零双边理想。
\end{theorem}
\begin{proof}
降链 \(A\supseteq J\supseteq J^2\supseteq\cdots\) 是有限维向量空间的降链，故存在 \(n\ge0\) 使 \(J^n=J^{n+1}\)。令 \(M=J^n\)，将它看成左 \(A\) 模。它有限维，因而有限生成；等式正是
\[
JM=J J^n=J^{n+1}=J^n=M.
\]
对理想 \(J\subseteq J(A)\) 和有限生成左模 \(M\) 使用\cref{b2:thm:noncommutative-nakayama}，得到 \(M=0\)。所以某个 \(J^n\) 为零，证明幂零性。

当 \(A\ne0\) 时，\(J\ne A\)，因为单位不属于任何真极大左理想。若在到达零以前某两个相邻幂相等，上面的同一论证会迫使它们为零，矛盾。因此从 \(A\) 到第一个零幂，每一步维数都严格下降，步数不超过 \(\dim_k A\)，得到所述界。最后，任意幂零双边理想由\cref{b2:prop:nilpotent-ideal-in-radical}包含于 \(J\)，而 \(J\) 本身已证幂零。
\end{proof}

证明用到了两个不同的事实：有限维性使某一项稳定，Nakayama 引理使这个稳定项为零。仅凭降链停止，只能得到 \(J^n=J^{n+1}\)，不能省去后一步。界 \(N\le\dim_k A\) 只把每个非零层的维数估为至少一，是一个统一的粗界，并不说最小幂零指数等于代数的维数。这个幂零性也将保证稍后提升幂等元的过程在有限步后结束。

% !TeX root = ../../Book2.tex
% 纲要标题：### 15.2 模的根、基座与 Loewy 列
% 纲要位置：计划纲要.md，第 1252 行。

\section{模的根、基座与 Loewy 列}
\label{b2:sec:1502}

现在把代数的根用于单个模。一个方向是不断取最大半单商，另一个方向是不断取出全部简单子模；它们给出从上到下和从下到上的两列。本节的 \(M\) 均为有限维左 \(A\) 模，\(J=J(A)\)。

% 纲要要点（供目录核对）：
%
% * 模的根与基座（socle；所有简单子模之和）
% * Loewy 列及半单层
% * 具体三角矩阵代数的计算
%

\subsection{模的根与基座}
\label{b2:subsec:1502:01}

\begin{definition}\label{b2:def:module-radical-socle}
设 \(M\) 为左 \(R\) 模。将 \(M\) 的全部极大真子模取交，所得子模称为 \(M\) 的\term[mo-de-gen]{根}[radical of a module]，记为 \(\operatorname{rad}M\)；若不存在极大真子模，该交按约定取 \(M\)。将全部简单子模求和，所得子模称为 \(M\) 的\term[jizuo]{基座}[socle]，记为 \(\operatorname{Soc}M\)；若没有简单子模，该和为零。对零模约定 \(\operatorname{rad}0=\operatorname{Soc}0=0\)。
\end{definition}
\symidx{\(\operatorname{rad}M\)、\(\operatorname{Soc}M\)：模的根与基座}

非零有限维模总有极大真子模及简单子模：分别取维数最大的真子模和维数最小的非零子模即可。基座是半单模，因为它是简单子模之和，适用\cref{b2:thm:semisimple-module-characterizations}。

\begin{theorem}\label{b2:thm:module-radical-jm}
设 \(k\) 为域、\(A\) 为有限维结合含幺 \(k\) 代数，\(J=J(A)\)，\(M\) 为有限维左 \(A\) 模，则
\[
\operatorname{rad}M=JM.
\]
商 \(M/JM\) 半单，而且对每个半单左 \(A\) 模 \(T\)，每个 \(A\) 模同态 \(M\to T\) 都唯一通过这个商分解。
\end{theorem}
\begin{proof}
对极大子模 \(L\)，商 \(M/L\) 简单，故被 \(J\) 零化。这说明 \(JM\subseteq L\)，对全部 \(L\) 取交得 \(JM\subseteq\operatorname{rad}M\)。

另一方面，\(M/JM\) 是有限维 \(A/J\) 模，由\cref{b2:prop:finite-algebra-semisimple-quotient}半单。若 \(m\notin JM\)，把它的非零商类投影到某个非零的简单直和分量，得到一个不将 \(m\) 送到零的满同态 \(M\to S\)。其核为不含 \(m\) 的极大子模，所以 \(m\notin\operatorname{rad}M\)。这给出反向包含。

最后，半单模是简单模的直和，所以 \(J\) 在其上作用为零。任意同态 \(f:M\to T\) 因而满足 \(f(JM)=J\cdot f(M)=0\)。由商模泛性质，它唯一通过 \(M/JM\) 分解。
\end{proof}

\begin{proposition}\label{b2:prop:module-socle-annihilator}
设 \(k\) 为域、\(A\) 为有限维结合含幺 \(k\) 代数，\(J=J(A)\)，\(M\) 为有限维左 \(A\) 模，则其基座为
\[
\operatorname{Soc}M=\{\,m\in M\mid Jm=0\,\}.
\]
\end{proposition}
\begin{proof}
每个简单子模均被 \(J\) 零化，所以基座包含于右侧。右侧是子模：若 \(Jm=0\)，则对 \(a\in A,j\in J\)，有 \(j(am)=(ja)m=0\)，因为 \(ja\in J\)。这个子模被 \(J\) 零化，由\cref{b2:prop:finite-algebra-semisimple-quotient}半单，故是若干简单子模之和，包含于基座。
\end{proof}

在 \(M=A=k[t]/(t^m)\)、\(m>1\) 中，\(\operatorname{rad}M=(\bar t)\)，而 \(\operatorname{Soc}M=k\bar t^{m-1}\)。前者收集不能通过简单商检测到的元素，后者收集已经构成简单子模的部分；它们通常不是同一个子空间。

\subsection{Loewy 列与半单层}
\label{b2:subsec:1502:02}

\begin{definition}\label{b2:def:loewy-filtrations}
设 \(M\) 为左 \(R\) 模。令 \(\operatorname{rad}^0M=M\)、\(\operatorname{rad}^{r+1}M=\operatorname{rad}(\operatorname{rad}^rM)\)（\(r\ge0\)），所得降列称为\term[gen-Loewy-lie]{根 Loewy 列}[radical Loewy series]。另令 \(\operatorname{Soc}^0M=0\)，并令 \(\operatorname{Soc}^{r+1}M\) 是 \(M/\operatorname{Soc}^rM\) 的基座在 \(M\) 中的原像，得到\term[jizuo-Loewy-lie]{基座 Loewy 列}[socle Loewy series]。
\end{definition}

以上两列对任意环上的模都有定义；下面用根的幂及其零化子写出的公式，适用于有限维代数上的有限维模。

\begin{proposition}\label{b2:prop:loewy-series}
设 \(k\) 为域、\(A\) 为有限维结合含幺 \(k\) 代数，\(J=J(A)\)，\(M\) 为有限维左 \(A\) 模。对所有 \(r\ge0\)，有
\[
\operatorname{rad}^rM=J^rM,\qquad
\operatorname{Soc}^rM=\{\,m\in M\mid J^rm=0\,\}.
\]
两列的相邻层均半单。若 \(J^N=0\)，根列至多在第 \(N\) 步到达零，基座列至多在第 \(N\) 步到达 \(M\)。
\end{proposition}
\begin{proof}
第一式由\cref{b2:thm:module-radical-jm}逐次用于有限维子模得到。第二式对 \(r\) 归纳；\(r=0\) 时 \(J^0=A\)，被 \(A\) 零化的元素只有零，符合定义。假设第 \(r\) 步成立，由\cref{b2:prop:module-socle-annihilator}，\(m\) 属于下一步当且仅当 \(Jm\subseteq\operatorname{Soc}^rM\)，即 \(J^rJm=0\)，也就是 \(J^{r+1}m=0\)。

根列的层 \(J^rM/J^{r+1}M\) 被 \(J\) 零化，故半单；基座列的层按定义是一个商模的基座，也半单。最后由\cref{b2:thm:finite-algebra-radical-nilpotent}取 \(J^N=0\)，代入两个公式便得到终止性。
\end{proof}

对有限维 \(k\) 代数 \(A\) 的有限维左模 \(M\)，记 \(J=J(A)\)。使 \(J^\ell M=0\) 的最小非负整数 \(\ell\) 称为 \(M\) 的\term[Loewy-changdu]{Loewy 长度}[Loewy length]。由命题，它也等于使 \(\operatorname{Soc}^\ell M=M\) 的最小整数。零模长度为零；非零半单模长度为一。

\ProseParagraphBreak
Loewy 层记录半单的商，不保证能在原模中为每一层同时选出子模代表。对双数代数 \(k[\varepsilon]/(\varepsilon^2)\) 的正则模，两层均为简单模 \(k\)，但层间扩张不分裂，因为若正则模为 \(k\oplus k\)，则 \(\varepsilon\) 应作用为零。这类层间关系正是稍后投射覆盖与箭图表示需要保存的信息。

\subsection{三角矩阵代数的计算}
\label{b2:subsec:1502:03}

令 \(A=T_3(k)\) 为上三角矩阵代数，矩阵单位记为 \(E_{ij}\)、\(1\le i\le j\le3\)。严格上三角矩阵组成理想 \(N\)，有 \(N^3=0\)，而 \(A/N\cong k^3\) 半单。由\cref{b2:prop:nilpotent-ideal-in-radical}与\cref{b2:prop:radical-quotient}，得到
\[
J=N=kE_{12}+kE_{13}+kE_{23},\qquad J^2=kE_{13},\qquad J^3=0.
\]
简单模 \(S_i\) 是一维空间，矩阵 \(a\) 通过对角元 \(a_{ii}\) 作用。它们给出全部简单模，因为简单模被 \(J\) 零化，因而来自 \(k^3\) 的一个坐标因子。

\ProseParagraphBreak
取 \(e_j=E_{jj}\)、\(P_j=Ae_j\)。它是第 \(j\) 列可能非零的矩阵组成的左理想，基为 \(E_{1j},\ldots,E_{jj}\)。矩阵单位乘法
\[
E_{ab}E_{ij}=\delta_{bi}E_{aj}
\]
给出
\begin{equation}\label{b2:eq:triangular-projective-layers}
J^rP_j=\operatorname{span}_k\{\,E_{ij}\mid1\le i\le j-r\,\},
\qquad 0\le r\le j,
\end{equation}
空指标集表示零。每个非零层由 \(E_{j-r,j}\) 的陪集生成；左乘矩阵 \(a\) 后，较低行号的项属于下一层，只剩系数 \(a_{j-r,j-r}\)，所以这个层同构于 \(S_{j-r}\)。

\ProseParagraphBreak
因此，从顶层向下读，\(P_1,P_2,P_3\) 的简单层依次为
\[
P_1:\ S_1;\qquad P_2:\ S_2,\ S_1;\qquad
P_3:\ S_3,\ S_2,\ S_1.
\]
它们的 Loewy 长度分别为一、二、三。另一方面，\(J\) 零化的向量恰为第一行号的部分，所以
\[
\operatorname{Soc}P_j=kE_{1j}\cong S_1,
\qquad
\operatorname{Soc}^rP_j=\operatorname{span}_k\{\,E_{ij}\mid i\le\min(r,j)\,\}.
\]
这些公式可直接从\cref{b2:prop:loewy-series}及矩阵乘法核验。

\ProseParagraphBreak
正则模分解为 \(A=P_1\oplus P_2\oplus P_3\)。于是
\[
A/JA\cong S_1\oplus S_2\oplus S_3,
\qquad \operatorname{Soc}A\cong S_1^{\oplus3}.
\]
最大半单商与最大半单子模甚至可以含有完全不同的重数。自然列向量模 \(k^3\) 则与 \(P_3\) 同构，映射为 \(E_{i3}\mapsto e_i\)；其根列是熟悉的坐标旗标，只是排列方向与从零开始的升旗标相反。

% !TeX root = ../../Book2.tex
% 纲要标题：### 15.3 投射覆盖
% 纲要位置：计划纲要.md，第 1541 行。

\section{投射覆盖}
\label{b2:sec:1503}

自由满射总存在，但其来源可能带有可以删去的投射分量。投射覆盖要求满射在模论意义下已经不能缩小来源。为了构造它，先要把半单商中的列模提升回来；这一步依赖幂等元提升，不能用商中的分解代替原代数中的分裂。本节的 \(A\) 仍有限维，\(J=J(A)\)；幂等元提升的两个定理则只需所写的环与幂零理想假设。

% 纲要要点（供目录核对）：
%
% * 极小生成与本质满射
% * 在构造投射覆盖前，就地证明幂零理想模的幂等元提升，并处理有限组正交幂等元的相容提升
% * 明确本质满射的模论含义为核是多余子模；从提升后的本原幂等元构造不可分解投射模及其到简单模的覆盖
% * 有限维代数中的投射覆盖；目标有限生成时来源的有限性，投射覆盖与最少自由生成的区别
% * 不可分解投射模与简单模
% * 由有限维性、半单商及根的幂零性逐项验证覆盖的存在和唯一性，避免把半单商中的分解直接当作原模中的分裂
% * 正文给出分别提升不保持正交的反例；唯一性包含目标同构时每个提升可逆，极小性说明投射目标的覆盖为同构
%

\subsection{极小生成与本质满射}
\label{b2:subsec:1503:01}

\begin{definition}\label{b2:def:superfluous-submodule}
设 \(R\) 为含幺环、\(P\) 为幺左 \(R\) 模、\(N\le P\) 为子模。若对每个子模 \(L\le P\)，等式 \(L+N=P\) 都蕴含 \(L=P\)，则称 \(N\) 为 \(P\) 的\term[duoyu-zimo]{多余子模}[superfluous submodule]，记为 \(N\ll P\)。若幺左 \(R\) 模的满同态 \(p:P\to M\) 满足 \(\ker p\ll P\)，则称 \(p\) 为\term[benzhi-manshe]{本质满射}[essential epimorphism]。
\end{definition}
\symidx{\(N\ll P\)：\(N\) 是 \(P\) 的多余子模}

这一条件也可写成：对每个子模 \(L\le P\)，若 \(p(L)=M\)，则 \(L=P\)。事实上，\(p(L)=M\) 等价于每个 \(x\in P\) 可写成 \(\ell+n\)，其中 \(\ell\in L,n\in\ker p\)，也就是 \(L+\ker p=P\)。因此任何仍能满射到目标的子模都已经是整个来源。多余性总是相对于指定的母模而言，判定时量化的是这个母模的全部子模。

\ProseParagraphBreak

等价地，对任意左 \(R\) 模同态 \(g:Q\to P\)，若复合 \(pg:Q\to M\) 满射，则 \(g\) 满射：此时 \(p(g(Q))=M\)，上一条件给出 \(g(Q)=P\)。反过来，把 \(g\) 取为任意子模 \(L\le P\) 的包含映射，就得到上一条件。

\begin{definition}\label{b2:def:projective-cover}
设 \(R\) 为含幺环、\(M\) 为幺左 \(R\) 模。若幺左 \(R\) 模的满同态 \(p:P\to M\) 的来源 \(P\) 投射，且 \(\ker p\ll P\)，则称 \(p\) 为 \(M\) 的一个\term[toushe-fugai]{投射覆盖}[projective cover]。定义中保留这个满同态，而不仅记录 \(P\) 的同构类型。
\end{definition}

若 \(M\) 有限生成，它的任意投射覆盖的来源也有限生成：取 \(M\) 的有限生成元，逐个提升到 \(P\)，令 \(L\) 为这些提升所生成的子模，则 \(p(L)=M\)。本质满射的条件给出 \(L=P\)。因此，在有限维代数 \(A\) 上，有限生成目标模的覆盖来源有限生成，进而有限维；这适用于本章讨论的有限维目标模。

\ProseParagraphBreak
\cref{b2:cor:nakayama-generators}说明，\(M/JM\) 的一组 \(A/J\) 模生成元可以提升为 \(M\) 的生成元。局部代数中，剩余模是在除环 \(A/J\) 上的向量空间，选基可得到最少的生成元。但一般情形下，由最少生成元给出的自由满射仍未必是投射覆盖。

\ProseParagraphBreak
例如，设 \(k\) 为域、\(A=M_2(k)\)、\(S=k^2\)。\(S\) 由一个非零向量生成，取第一列的映射 \(A\to S\) 因而用了最少的一个模生成元。不过真左理想 \(AE_{11}\) 已经满射到 \(S\)，故这个自由满射不是本质满射。限制 \(AE_{11}\to S\) 则是同构；\(AE_{11}\) 是正则模的直和因子，因而投射，这才给出一个投射覆盖。\(AE_{11}\) 的 \(k\) 维数为二，不能是 \(A\) 上的自由模。因此，最少生成元使自由来源的秩最小，来源仍可能缩小为一个非自由的投射模。

\subsection{幂等元与正交幂等元组的提升}
\label{b2:subsec:1503:02}

\begin{theorem}\label{b2:thm:idempotent-lifting}
设 \(R\) 为含幺环，\(I\) 为幂零双边理想。每个幂等元 \(\bar e\in R/I\) 都能提升为幂等元 \(e\in R\)。此外，两个这样的提升 \(e,f\) 共轭：存在 \(u\in1+I\)，使 \(f=ueu^{-1}\)。幂零假设不能一般省去：在 \(R=\mathbb Z\)、\(I=6\mathbb Z\) 中，幂等元 \(3+6\mathbb Z\) 不能提升为整数幂等元。
\end{theorem}
\begin{proof}
先处理 \(I^2=0\)。任取 \(a\) 提升 \(\bar e\)，令 \(d=a^2-a\in I\)，这是 \(a\) 未必幂等的误差。希望取 \(e=a+c\)，其中 \(c\in I\) 且与 \(a\) 交换。因为 \(I^2=0\) 使 \(c^2=0\)，所需条件化为
\[
e^2-e=d+ac+ca-c=d+(2a-1)c=0.
\]
这个误差修正方程对 \(c\) 是线性的。又有 \((2a-1)^2=1+4d\)；对任意 \(x\in I\)，\(dx\in I^2=0\)，所以 \((2a-1)^2x=x\)。因此在 \(I\) 上左乘 \(2a-1\) 的算子平方为恒等，可以取
\[
c=-(2a-1)d=(1-2a)d.
\]
所取 \(d,c\) 都是 \(a\) 的多项式，故确实与 \(a\) 交换；它们也都在 \(I\) 中。代入得到
\[
e^2-e=d+(2a-1)c
=d-(2a-1)^2d.
\]
由于 \((2a-1)^2=1+4d\)，右侧等于 \(-4d^2=0\)。因此 \(e\) 幂等，且与 \(a\) 模 \(I\) 相同。这里没有除以二，对任意特征都成立。

一般地，取 \(I^N=0\)。已在 \(R/I^r\) 中得到的幂等元，可以提升到 \(R/I^{r+1}\)：自然满射的核为 \(I^r/I^{r+1}\)，它的平方为零，因为 \(2r\ge r+1\)、\(r\ge1\)。将平方零情形用于这个商环，依次作 \(r=1,\ldots,N-1\) 的提升，最后得到 \(R/I^N=R\) 中的幂等元。

为证共轭性，先把 \(e,f\) 看成正则右模 \(R_R\) 上的左乘投影；左乘是右 \(R\) 线性的，因为 \((ex)r=e(xr)\)。它们分别给出分解
\[
R_R=eR\oplus(1-e)R=fR\oplus(1-f)R.
\]
我们寻找的左乘映射应把第一组分解的两个部分分别送到第二组对应的部分，从而满足 \(ue=fu\)。对 \(eR\) 中的元素左乘 \(f\)，对 \((1-e)R\) 中的元素左乘 \(1-f\)，便得到候选元素
\[
u=fe+(1-f)(1-e).
\]
事实上，对 \(x\in eR\) 有 \(ux=fx\in fR\)，对 \(x\in(1-e)R\) 有 \(ux=(1-f)x\in(1-f)R\)；直接相乘也得到 \(ue=fe=fu\)。这个公式先完成分解的匹配，还须保证匹配映射可逆。因 \(\bar f=\bar e\)，模 \(I\) 后有 \(\bar u=\bar e+(1-\bar e)=1\)，故 \(u\in1+I\)。幂零性使 \(u\) 可逆，可由有限几何和写出其逆。于是 \(ue=fu\) 给出 \(f=ueu^{-1}\)。

最后，\(3^2-3=6\)，所以 \(3+6\mathbb Z\) 幂等。整数幂等元满足 \(e(e-1)=0\)，而 \(\mathbb Z\) 是整环，故只有零或一；它们模六都不等于三，说明这个幂等元不能提升。这里 \((6\mathbb Z)^r=6^r\mathbb Z\ne0\) 对每个正整数 \(r\) 成立，恰好缺少幂零假设。
\end{proof}

要提升一组分解，所选提升还须保持彼此正交和总和为一；分别任意提升每个幂等元并不足够。\cref{b2:prop:idempotent-lift-compatibility-counterexample}在本节末给出各自幂等而不相容的完整计算。下一定理把这种相容性一并保证。

\begin{theorem}\label{b2:thm:orthogonal-idempotent-lifting}
设 \(R\) 为含幺环、\(I\subseteq R\) 为幂零双边理想，\(m\) 为正整数。若 \(\bar e_1,\ldots,\bar e_m\in R/I\) 满足
\[
\bar e_i\bar e_j=\delta_{ij}\bar e_i,
\qquad \sum_{i=1}^m\bar e_i=1,
\]
则存在提升 \(e_1,\ldots,e_m\in R\)，同样满足 \(e_ie_j=\delta_{ij}e_i\) 与 \(\sum_i e_i=1\)。
\end{theorem}
\begin{proof}
对 \(m\) 归纳。\(m=1\) 时取 \(e_1=1\)。设 \(m>1\)，先用\cref{b2:thm:idempotent-lifting}提升 \(\bar e_1\)，并令 \(c=1-e_1\)。余下部分的单位应为 \(c\)，所以转到角环 \(cRc\) 中提升其余幂等元。映射
\[
cRc\longrightarrow\bar c(R/I)\bar c
\]
满射：右侧元素可由任意代表元两边乘 \(c\) 得到。其核为 \(cIc\)，因为同时属于 \(I\) 与 \(cRc\) 的元素 \(x\) 满足 \(x=cxc\)。这个核幂零，\((cIc)^N\subseteq cI^Nc=0\)。

余下的 \(\bar e_2,\ldots,\bar e_m\) 都在该商角环中，仍相互正交，总和为 \(\bar c\)。在 \(cRc\) 中应用归纳假设，得到正交提升 \(e_2,\ldots,e_m\)，总和为 \(c\)。它们均被 \(e_1\) 在左右两侧乘为零，所以连同 \(e_1\) 得到所求完整正交组。
\end{proof}

平方零修正与角环归纳可参阅 Etingof 等人的《Introduction to Representation Theory》\cite[Proposition 9.1.1; Corollary 9.1.3, pp. 209--210]{EtingofEtAl2011RepresentationTheory}。本书已把单个提升和整组相容提升分别展开；后面的构造将同时使用这两个结论。

\subsection{多余子模与简单模的投射覆盖}
\label{b2:subsec:1503:03}

\begin{proposition}\label{b2:prop:finite-superfluous-radical}
设 \(k\) 为域、\(A\) 为有限维结合含幺 \(k\) 代数，\(J=J(A)\)。有限维幺左 \(A\) 模 \(P\) 的子模 \(N\) 多余，当且仅当 \(N\subseteq JP=\operatorname{rad}P\)。因此，对有限生成投射幺左 \(A\) 模的满射 \(p:P\to M\)，以下条件等价：它是投射覆盖；\(\ker p\subseteq JP\)；诱导映射 \(P/JP\to M/JM\) 是同构。若这个满射是投射覆盖且目标 \(M\) 本身投射，则 \(p\) 为同构。
\end{proposition}
\begin{proof}
若 \(N\ll P\)，而某个极大子模 \(L\) 不包含 \(N\)，则 \(L\subsetneq L+N\)，极大性迫使 \(L+N=P\)，与 \(L\ne P\) 矛盾。因此 \(N\) 包含于所有极大子模的交，由\cref{b2:thm:module-radical-jm}得到 \(N\subseteq JP=\operatorname{rad}P\)。

反过来，若 \(N\subseteq JP\) 且 \(L+N=P\)，则 \(L+JP=P\)。商 \(P/L\) 有限生成，并满足 \(J(P/L)=P/L\)，由\cref{b2:thm:noncommutative-nakayama}，\(P/L=0\)，即 \(L=P\)。

对满射 \(p\)，有 \(p(JP)=JM\)。诱导的商映射满射，核为 \((\ker p+JP)/JP\)：若 \(p(x)\in JM\)，取 \(y\in JP\) 使 \(p(y)=p(x)\)，便有 \(x-y\in\ker p\)。所以该核为零恰等于 \(\ker p\subseteq JP\)，得到其余等价性。

若 \(p\) 是覆盖且 \(M\) 投射，取截面 \(s:M\to P\)，便有 \(P=\ker p\oplus s(M)\)。核多余使 \(s(M)=P\)，从而 \(\ker p=0\)，\(p\) 是同构。
\end{proof}

\begin{definition}\label{b2:def:primitive-idempotent}
设 \(k\) 为域、\(A\) 为结合含幺 \(k\) 代数、\(e\in A\) 为非零幂等元。若 \(e\) 不能写成两个非零正交幂等元之和，则称它为\term[benyuan-midengyuan]{本原幂等元}[primitive idempotent]。这里的正交要求两个方向的乘积都为零。
\end{definition}

这里“本原”修饰幂等元，表示它不能再作正交分解；\cref{b2:def:primitive-semiprimitive}中的本原环则要求环有忠实简单模，两者的定义不同。

\begin{proposition}\label{b2:prop:primitive-projective-indecomposable}
设 \(k\) 为域、\(A\) 为结合含幺 \(k\) 代数、\(e\in A\) 为非零幂等元。幺左 \(A\) 模 \(Ae\) 有限生成且投射；它不可分解，当且仅当 \(e\) 本原。
\end{proposition}
\begin{proof}
\(Ae\) 由单个元素 \(e\) 生成。正则左模具有分解 \(A=Ae\oplus A(1-e)\)，投影为右乘 \(e\)，所以 \(Ae\) 是有限自由模的直和因子，因而投射。若 \(e=f+g\) 是非零正交幂等元之和，则 \(Ae=Af\oplus Ag\)，给出非平凡分解。

反过来，若 \(Ae\) 有非平凡分解，对应投影是 \(\operatorname{End}_A(Ae)\) 中一个非零且非恒等的幂等元。由\cref{b2:prop:corner-endomorphism}，该自同态环同构于 \((eAe)^{\mathrm{op}}\)，所以角环中也有 \(f^2=f\)、\(f\ne0,e\)。于是 \(e=f+(e-f)\)，且 \(f(e-f)=(e-f)f=0\)，这与本原性矛盾。
\end{proof}

\begin{proposition}\label{b2:prop:primitive-projective-cover}
设 \(k\) 为域、\(A\) 为有限维结合含幺 \(k\) 代数，\(J=J(A)\)。取 \(A/J\) 的本原幂等元 \(\bar e\)，并取其在 \(A\) 中的幂等元提升 \(e\)，则 \(e\) 本原，\(Ae/Je\) 是简单幺左 \(A\) 模，且自然满射
\[
Ae\longrightarrow Ae/Je
\]
是投射覆盖。每个简单幺左 \(A\) 模都能由这样的满射得到覆盖。
\end{proposition}
\begin{proof}
先看单个矩阵环 \(M_n(D)\)，其中 \(D\) 为除环。幂等矩阵作为右 \(D\) 向量空间上的算子，其像与核给出直和；选择相应基后，它共轭于 \(\operatorname{diag}(I_r,0)\)。后者是 \(r\) 个秩一坐标投影之和，所以 \(r>1\) 时不本原。若 \(r=1\)，两个非零正交幂等元的像会给出两个非零直和分量，不能同时放在这个一维像中。因此非零幂等元本原当且仅当 \(r=1\)。再看\eqref{b2:eq:finite-algebra-semisimple-quotient}中的有限矩阵环直积，本原幂等元只能在一个矩阵因子中具有秩一，其余因子为零。相应左理想同构于一个列模，且简单：若列 \(v\) 的第 \(j\) 个坐标非零，为把它映到任意给定列 \(w\)，取矩阵 \(C\) 的第 \(j\) 列为 \((w_iv_j^{-1})_i\)，其余列为零，便有 \(Cv=w\)。

\(J\) 零化每个简单左 \(A\) 模，所以它们都可看作 \(A/J\) 的简单模。半单商的正则模是上述简单列模的有限直和；任意简单模都是这个正则模的商。相应满射在至少一个列模上的限制非零，由\cref{b2:lem:schur}，该限制是同构。这也说明上述构造遍及全部简单模。

若 \(e=f+g\) 是正交幂等元分解，模 \(J\) 后，\(\bar e\) 的本原性使 \(\bar f,\bar g\) 至少一个为零。该幂等元便属于幂零理想 \(J\)，只能为零，因为幂等元的每个正幂都等于自身。因此 \(e\) 本原。最后，\(Ae\to(A/J)\bar e\) 的核为 \(Je\)：核元素既在 \(J\) 中，又满足 \(x=xe\)，故在 \(Je\) 中。于是 \(Ae/Je\) 是刚才的简单列模，且核为 \(J(Ae)\)。由\cref{b2:prop:finite-superfluous-radical}，满射为投射覆盖。
\end{proof}

\subsection{投射覆盖的存在性}
\label{b2:subsec:1503:04}

从 \(A/J\) 的每个矩阵因子选出全部对角矩阵单位，利用\cref{b2:thm:orthogonal-idempotent-lifting}得到 \(A\) 中的完整正交提升。每个提升给出上述简单列模的覆盖；同一矩阵因子的列模同构，但暂不预设其提升模已经同构。

\ProseParagraphBreak
从每一种简单模的同构类型中取一个代表 \(S_1,\ldots,S_s\)，并固定由\cref{b2:prop:primitive-projective-cover}构造的覆盖 \(p_i:P_i\to S_i\)。于是 \(P_i/JP_i\cong S_i\)，且 \(P_i\) 不可分解。

\begin{theorem}\label{b2:thm:projective-cover-existence}
设 \(k\) 为域、\(A\) 为有限维结合含幺 \(k\) 代数，\(J=J(A)\)。从全部简单幺左 \(A\) 模的同构类中各选一个代表 \(S_1,\ldots,S_s\)，并给每个 \(S_i\) 选定投射覆盖 \(p_i:P_i\to S_i\)，则每个有限维幺左 \(A\) 模 \(M\) 都有投射覆盖。若
\[
M/JM\cong\bigoplus_{i=1}^s S_i^{\oplus m_i},
\]
可以取覆盖的来源为 \(P=\bigoplus_iP_i^{\oplus m_i}\)。
\end{theorem}
\begin{proof}
由\cref{b2:thm:module-radical-jm}，所写的有限半单分解存在。合并 \(p_i\) 的相应份数并接上这个同构，得到满同态 \(h:P\to M/JM\)，其在商上诱导同构 \(P/JP\cong M/JM\)。\(P\) 是有限个投射模的直和，因而投射，可以沿商映射 \(\pi:M\to M/JM\) 提升 \(h\)，得到 \(p:P\to M\) 且 \(\pi p=h\)。

因为 \(h\) 满射，有 \(p(P)+JM=M\)。对有限生成的 \(M\) 及其子模 \(p(P)\) 应用\cref{b2:cor:nakayama-generators}的生成元提升形式，得到 \(p(P)=M\)。又因为 \(p\) 在模 \(J\) 后是同构，\(\ker p\subseteq JP\)。由\cref{b2:prop:finite-superfluous-radical}，这个满射正是投射覆盖。
\end{proof}

这里并未把 \(M/JM\) 的简单直和分解直接提升成 \(M\) 的简单直和分解。被提升的是它们的投射覆盖，以及一个从投射模出发的映射；满射性由 Nakayama 引理重新保证。

\ProseParagraphBreak
例如在 \(A=k[t]/(t^3)\) 中，\(M=A/(\bar t^2)\) 的最大半单商为 \(k\)，对应投射覆盖是 \(A\to M\)。核 \((\bar t^2)\) 包含在 \(JA\) 中，所以多余；这个覆盖保留了 \(M\) 的两层结构，而来源 \(A\) 本身有三层。

\subsection{不可分解投射模与简单模}
\label{b2:subsec:1503:05}

\begin{theorem}\label{b2:thm:indecomposable-projective-simple}
设 \(k\) 为域、\(A\) 为有限维结合含幺 \(k\) 代数，\(J=J(A)\)。从全部简单幺左 \(A\) 模的同构类中各选一个代表 \(S_1,\ldots,S_s\)，并按\cref{b2:prop:primitive-projective-cover}给每个 \(S_i\) 选定形如 \(p_i:P_i=Ae_i\to S_i\) 的投射覆盖。于是每个有限生成投射幺左 \(A\) 模都是 \(P_1,\ldots,P_s\) 的有限直和。有限生成不可分解投射模恰为这些 \(P_i\) 的同构类型，而
\[
P\longmapsto P/JP
\]
给出有限生成不可分解投射模与简单模的同构类型之间的一一对应。
\end{theorem}
\begin{proof}
对有限生成投射模 \(Q\)，按\cref{b2:thm:projective-cover-existence}构造覆盖 \(p:P\to Q\)，其中 \(P\) 是若干 \(P_i\) 的直和。\(Q\) 投射，而 \(p\) 是投射覆盖，由\cref{b2:prop:finite-superfluous-radical}，\(p\) 本身就是同构，故 \(Q\cong P\)。

各 \(P_i\) 已由\cref{b2:prop:primitive-projective-cover}证明不可分解；若 \(Q\) 不可分解，上述直和只能有一个非零分量。不同的 \(P_i\) 不同构，因为其商 \(P_i/JP_i\cong S_i\) 不同构。每个简单模都有其中一个覆盖，故得到双射。
\end{proof}

因此，这里对应的是三种对象的同构类型：
\[
S_i\quad\longleftrightarrow\quad P_i\cong Ae_i
\quad\longleftrightarrow\quad
\{\,e\text{ 为本原幂等元}\mid Ae\cong Ae_i\,\}.
\]
右端是一类本原幂等元，而不是某一个唯一的幂等元。

\ProseParagraphBreak
直和中的重数也可以直接从最大半单商读出。对有限生成投射左 \(A\) 模 \(Q\)，令 \(\Delta_i=\operatorname{End}_A(S_i)\)，它由\cref{b2:lem:schur}为除环。\(\operatorname{Hom}_A(Q,S_i)\) 通过像端的复合成为左 \(\Delta_i\) 向量空间，具体为 \(\alpha\cdot f=\alpha\circ f\)。因为 \(J\) 零化 \(S_i\)，每个同态都通过 \(Q/JQ\) 分解，因而
\begin{equation}\label{b2:eq:projective-simple-multiplicity}
Q\cong\bigoplus_iP_i^{\oplus m_i}
\quad\Longrightarrow\quad
m_i=\dim_{\Delta_i}\operatorname{Hom}_A(Q,S_i).
\end{equation}
理由是不同简单模之间没有非零同态，\(\operatorname{Hom}_A(S_i,S_i)=\Delta_i\) 在自身上的左维数为一。这也证明这些重数唯一，且没有假设 \(\Delta_i=k\)。\cref{b2:prop:semisimple-multiplicity}从来源简单模的自同态作预复合，得到右除环向量空间；这里简单模位于像端，作后复合，因而得到左除环向量空间。

\subsection{投射覆盖的唯一性与极小性}
\label{b2:subsec:1503:06}

存在性已经由\cref{b2:thm:projective-cover-existence}得到。以下唯一性须比较连同映射的覆盖；只比较来源的维数，不足以说明它们在 \(M\) 上相容。

\begin{theorem}\label{b2:thm:projective-cover-uniqueness}
设 \(k\) 为域、\(A\) 为有限维结合含幺 \(k\) 代数，\(M,N\) 为有限维幺左 \(A\) 模，\(h:M\to N\) 为 \(A\) 模同构，\(p:P\to M\) 与 \(q:Q\to N\) 为投射覆盖。则存在满足 \(qa=hp\) 的 \(A\) 模同态 \(a:P\to Q\)，且每个满足这一等式的 \(a\) 都是同构。取 \(N=M\)、\(h=1_M\)，便得到两个覆盖在目标 \(M\) 上同构；这样的同构本身未必唯一。
\end{theorem}
\begin{proof}
由 \(P\) 的投射性，可以沿满射 \(q\) 提升 \(hp\)，所以满足 \(qa=hp\) 的 \(a\) 存在。现在固定任意一个这样的 \(a\)。由 \(Q\) 的投射性，沿 \(p\) 提升 \(h^{-1}q\)，得到 \(b:Q\to P\)，满足 \(pb=h^{-1}q\)。于是 \(pba=p\)、\(qab=q\)。对任意 \(x\in P\)，有 \(x-ba(x)\in\ker p\)，所以
\[
ba(P)+\ker p=P.
\]
核多余，故 \(ba(P)=P\)。同样由 \(qab=q\) 及 \(\ker q\ll Q\) 得 \(ab(Q)=Q\)。这两步只用了覆盖的多余核条件。

目标 \(M,N\) 有限生成，本质满射性使来源 \(P,Q\) 也有限生成；又因 \(A\) 有限维，\(P,Q\) 都有限维，因而具有有限长度。\cref{b2:prop:finite-length-endomorphism}使满自同态 \(ba\) 与 \(ab\) 可逆。于是 \((ba)^{-1}b\) 是 \(a\) 的左逆，而 \(b(ab)^{-1}\) 是其右逆；一个映射的左逆与右逆相同，故给定的 \(a\) 是同构。

例如 \(A=k[t]/(t^2)\) 的覆盖 \(A\to k\) 上，恒等映射和乘以 \(1+\bar t\) 的映射都是同构，且模 \(\bar t\) 后都诱导恒等，二者不同。因此唯一的是覆盖在目标上同构的类型。
\end{proof}

\begin{proposition}\label{b2:prop:projective-cover-minimality}
设 \(k\) 为域、\(A\) 为有限维结合含幺 \(k\) 代数，\(M\) 为幺左 \(A\) 模。对有限生成投射幺左 \(A\) 模的满射 \(p:P\to M\)，它是投射覆盖，当且仅当每个满足 \(pu=p\) 的 \(A\) 模自同态 \(u:P\to P\) 都可逆。若 \(p\) 是覆盖，且 \(q:Q\to M\) 是任意投射幺左 \(A\) 模的满射，则 \(P\) 是 \(Q\) 的直和因子，不要求 \(Q\) 有限生成。
\end{proposition}
\begin{proof}
覆盖情形中，\(u(P)+\ker p=P\)，所以 \(u\) 满射。由命题假设，\(P\) 在有限维 \(k\) 代数 \(A\) 上有限生成，故它本身有限维，\(u\) 因而可逆。反过来，若 \(L+\ker p=P\)，则 \(p|_L:L\to M\) 满射。利用 \(P\) 的投射性，把 \(p\) 提升到 \(L\)，再复合包含映射得到 \(u:P\to P\)，其像在 \(L\) 中且 \(pu=p\)。若它必可逆，则 \(L=P\)，所以核多余。

对第二个断言，投射性给出 \(a:P\to Q\)、\(b:Q\to P\)，满足 \(qa=p\)、\(pb=q\)。于是 \(pba=p\)，第一部分使 \(ba\) 可逆。映射 \(a(ba)^{-1}:P\to Q\) 是 \(b\) 的截面，因此 \(Q\cong P\oplus\ker b\)。
\end{proof}

这使“覆盖极小”有了可检验的含义：任何其他投射满射的来源都含有它作为直和因子；若一个来源还能通过自同态缩小却仍覆盖目标，就不满足这种极小性。半单代数上的有限维模都投射，所以它们的覆盖不会增添额外的层，来源经覆盖映射与目标同构。关于同构类的选择虽然唯一，关于具体提升映射的选择仍须在后续构造中保留。

\begin{proposition}\label{b2:prop:idempotent-lift-compatibility-counterexample}
设 \(k\) 为域、\(R=M_2\bigl(k[\varepsilon]/(\varepsilon^2)\bigr)\)，\(\varepsilon\) 的商类仍记为 \(\varepsilon\)，\(E_{ij}\) 为矩阵单位，\(I=\varepsilon R\)。虽然 \(I^2=0\)，幂等元
\[
e=E_{11}+\varepsilon E_{12},\qquad f=E_{22}
\]
分别提升了 \(R/I\) 中正交且和为一的 \(\bar E_{11},\bar E_{22}\)，所选 \(e,f\) 却不正交。若改取 \(f'=1-e=E_{22}-\varepsilon E_{12}\)，则 \(e,f'\) 为同一商幂等元组的正交提升，且和为一。
\end{proposition}
\begin{proof}
\(I\) 为双边理想，且 \(\varepsilon^2=0\) 给出 \(I^2=0\)。由矩阵单位乘法，\(E_{11}E_{12}=E_{12}\)、\(E_{12}E_{11}=0\)，所以
\[
e^2=E_{11}+\varepsilon E_{12}=e,\qquad f^2=f.
\]
模 \(I\) 后，它们分别为 \(E_{11},E_{22}\)，其乘积在两个方向都为零，和为单位。但是在 \(R\) 中，
\[
ef=\varepsilon E_{12}\ne0,\qquad fe=0.
\]
因此分别取出的幂等提升不相互正交，它们的和也为 \(1+\varepsilon E_{12}\)，而不是一。取 \(f'=1-e\)，由 \(e^2=e\) 直接得到 \((f')^2=f'\)、\(ef'=f'e=0\) 及 \(e+f'=1\)；同时 \(f'-f\in I\)，所以没有改变第二个商幂等元。所有计算均适用于任意特征。
\end{proof}

% !TeX root = ../../Book2.tex
% 纲要标题：### 15.4 基本代数的动机
% 纲要位置：计划纲要.md，第 1267 行。

\section{基本代数的动机}
\label{b2:sec:1504}

正则模中同一种不可分解投射模可能出现多次；矩阵大小的一部分信息正来自这种重复。本节从每一种同构类型中只保留一个，构造一个较小的角代数，并说明它为什么适合记录简单模之间的关系。本节只证明这个角代数的基本性质，模范畴的等价留在下一章建立。

% 纲要要点（供目录核对）：
%
% * 本原幂等元
% * 简单模的自同态除环
% * 基本与分裂基本代数；第五卷箭图表示的代数来源
%
% ---
%

\subsection{本原幂等元}
\label{b2:subsec:1504:01}

\cref{b2:def:primitive-idempotent}与\cref{b2:prop:primitive-projective-indecomposable}已经把本原幂等元与不可分解投射左理想联系起来。本原性不要求幂等元位于中心：例如上三角矩阵代数的 \(E_{ii}\) 本原，却通常不是中心元。正交分解一的用途，是分解左正则模；只有各项还位于中心时，才能把它解释成代数的直积。

\begin{proposition}\label{b2:prop:corner-radical}
设 \(k\) 为域、\(A\) 为有限维结合含幺 \(k\) 代数，\(J=J(A)\)。对幂等元 \(e\in A\)，记 \(\bar e=e+J\in A/J\)。角代数 \(B=eAe\) 满足
\[
J(B)=eJe,
\qquad
B/J(B)\cong\bar e(A/J)\bar e.
\]
若 \(e\ne0\)，则 \(e\) 本原当且仅当 \(\bar e\) 本原；此时右侧是除环，\(eAe\) 是局部环。
\end{proposition}
\begin{proof}
向商代数映射并取角给出满射 \(eAe\to\bar e(A/J)\bar e\)，核为 \(eJe\)，这与\cref{b2:thm:orthogonal-idempotent-lifting}证明中的角环商计算相同。\(eJe\) 幂零，因为其 \(N\) 次乘积包含在 \(eJ^Ne\) 中，而 \(J^N=0\)。因此 \(eJe\subseteq J(B)\)。

在\eqref{b2:eq:finite-algebra-semisimple-quotient}的每个矩阵块中，对幂等算子 \(\bar e\) 选择像与核的基，使它成为 \(\operatorname{diag}(I_{r_i},0)\)。相应角环为 \(M_{r_i}(D_i)\)，零秩的块省略。所以商角环是半单代数，根为零。由\cref{b2:prop:radical-quotient}，\(J(B)/eJe=0\)，故 \(J(B)=eJe\)。

若 \(\bar e\) 本原，\cref{b2:prop:primitive-projective-cover}已证 \(e\) 本原。反过来，若 \(e\) 本原，则 \(Ae\) 不可分解，由\cref{b2:thm:indecomposable-projective-simple}，\(Ae/Je\cong(A/J)\bar e\) 简单。上述矩阵块描述迫使恰有一个 \(r_i=1\)，其余为零，故 \(\bar e\) 本原，商角环是除环。

最后，在含幺环 \(B\) 中，模 \(J(B)\) 后可逆的元素本身可逆。确实，若 \(xy\) 与 \(yx\) 都模 \(J(B)\) 等于单位，则它们由\cref{b2:thm:jacobson-unit-criterion}可逆，从而 \(x\) 同时有右逆 \(y(xy)^{-1}\) 与左逆 \((yx)^{-1}y\)，两者相同。于是当 \(B/J(B)\) 是除环时，\(B\) 的非单位元恰为 \(J(B)\)，它是唯一极大左理想，故 \(B\) 局部。此处角环的单位是 \(e\)。
\end{proof}

本原性在这里能传过商，依赖有限维情形下 \(J\) 的幂零性；对任意理想取商，幂等元的提升与上述结论都不能直接沿用。

\ProseParagraphBreak

例如 \(A=k[t]/(t^m)\) 的单位 \(1\) 是本原幂等元：该代数局部，模根后只有一个除环 \(k\)。当 \(m>1\) 时，角代数就是 \(A\) 自身，它虽局部却不是除环。局部性允许存在非零根，这与简单商是除环并不矛盾。

\subsection{简单模的自同态除环}
\label{b2:subsec:1504:02}

固定简单模 \(S_i\) 及其覆盖 \(P_i\)，仍记 \(\Delta_i=\operatorname{End}_A(S_i)\)。Schur 引理使 \(\Delta_i\) 成为除环，而不是预先成为基域。它控制\eqref{b2:eq:projective-simple-multiplicity}中的重数计数。

\begin{proposition}\label{b2:prop:projective-endomorphism-residue}
设 \(k\) 为域、\(A\) 为有限维结合含幺 \(k\) 代数，\(J=J(A)\)。令 \(S_i\) 为简单左 \(A\) 模，\(p_i:P_i\to S_i\) 为其投射覆盖，并记 \(\Delta_i=\operatorname{End}_A(S_i)\)。沿 \(P_i/JP_i\cong S_i\) 取商，从 \(P_i\) 的自同态得到满同态
\[
\operatorname{End}_A(P_i)\longrightarrow\operatorname{End}_A(S_i)=\Delta_i.
\]
其核是把 \(P_i\) 映入 \(JP_i\) 的全部自同态，恰为 \(\operatorname{End}_A(P_i)\) 的 Jacobson 根。
\end{proposition}
\begin{proof}
每个自同态保持 \(JP_i\)，所以诱导商上的自同态。若给定 \(\alpha:S_i\to S_i\)，利用 \(P_i\) 的投射性，把 \(\alpha p_i:P_i\to S_i\) 沿 \(p_i:P_i\to S_i\) 提升，就得到所需原像，证明满射。

核的描述来自商映射的定义。记此核为 \(K\)。若 \(u\in K\)，则对每个 \(r\) 有 \(u(J^rP_i)\subseteq J^r\cdot u(P_i)\subseteq J^{r+1}P_i\)。因而 \(N\) 个核中自同态的复合把 \(P_i\) 映入 \(J^NP_i=0\)，所以 \(K\) 是幂零理想。它包含于该自同态环的根，而模掉它的商是除环 \(\Delta_i\)，根为零；再用\cref{b2:prop:radical-quotient}，得到核恰为根。
\end{proof}

若 \(P_i=Ae\)，由\cref{b2:prop:corner-endomorphism}，这个自同态环是 \((eAe)^{\mathrm{op}}\)。因此半单商矩阵块里的除环与简单模自同态除环相差一个反环；按本书的左模约定，可在\eqref{b2:eq:finite-algebra-semisimple-quotient}中选取
\[
\Delta_i=\operatorname{End}_A(S_i),\qquad D_i=\Delta_i^{\mathrm{op}}.
\]
非交换除环上不能省略这个反环。

\begin{proposition}\label{b2:prop:split-simple-endomorphisms}
设 \(k\) 为代数闭域、\(A\) 为有限维结合含幺 \(k\) 代数。每个有限维简单左 \(A\) 模 \(S\) 都满足 \(\operatorname{End}_A(S)=k\operatorname{id}_S\)。
\end{proposition}
\begin{proof}
这是\cref{b2:cor:schur-algebraically-closed}的应用：本章的简单模有限维，基域代数闭，故每个简单自同态都是标量。该推论不要求 \(A\) 本身半单，因此可以直接用于这里的有限维代数。
\end{proof}

在一般基域上，此结论不成立。例如把 \(\mathbb C\) 看成实代数，简单左模 \(S=\mathbb C\) 的自同态除环是 \(\mathbb C\)，并不等于 \(\mathbb R\)。相应的 Hom 空间在基域上的维数为二，在这个自同态除环上的维数才为一。

\subsection{箭图表示的代数来源}
\label{b2:subsec:1504:03}

\begin{definition}\label{b2:def:basic-algebra}
设 \(k\) 为域、\(B\) 为有限维结合含幺 \(k\) 代数。若 \(B/J(B)\) 是有限个除 \(k\) 代数的直积，即其 Artin--Wedderburn 分解的每个矩阵块大小都为一，则称 \(B\) 为\term[jiben-daishu]{基本代数}[basic algebra]。
若进一步有 \(B/J(B)\cong k^s\)，则称 \(B\) 为\term[fenlie-jiben-daishu]{分裂基本代数}[split basic algebra]。
\end{definition}

\begin{proposition}\label{b2:prop:basic-corner}
设 \(k\) 为域、\(A\) 为有限维结合含幺 \(k\) 代数。可以选出幂等元 \(e\in A\)，使 \(eAe\) 为基本 \(k\) 代数，并满足 \(AeA=A\)。具体地，可在半单商 \(A/J(A)\) 的每个矩阵因子中只保留一个对角本原幂等元，再把这些幂等元相容提升并求和。
\end{proposition}
\begin{proof}
把\eqref{b2:eq:finite-algebra-semisimple-quotient}中全部对角矩阵单位组成完整正交组，由\cref{b2:thm:orthogonal-idempotent-lifting}提升为 \(e_{ij}\)，其中 \(i\) 标记矩阵因子，\(1\le j\le n_i\)。令 \(e=\sum_i e_{i1}\)，正交性保证它幂等。由\cref{b2:prop:corner-radical}，
\[
eAe/J(eAe)\cong\bar e(A/J)\bar e\cong\prod_iD_i,
\]
所以这个角代数基本。

在每个矩阵块中，一个对角矩阵单位生成整个双边理想：\(E_{p1}E_{11}E_{1q}=E_{pq}\)。因此 \(\bar e\) 在 \(A/J\) 中生成全部商代数，得到 \(AeA+J=A\)。写 \(1=x+j\)，其中 \(x\in AeA,j\in J\)。由根的可逆性判据，\(x=1-j\) 可逆；包含单位元的双边理想只能是整个 \(A\)，故 \(AeA=A\)。
\end{proof}

构造中同时得到 \(AeA=A\)：只保留每种列模的一份，仍能由角上的元素生成整个代数。这正是下一章比较模范畴所需的条件。设 \(A\) 为含幺代数、\(e\in A\) 为幂等元。若 \(AeA=A\)，则称 \(e\) 为\term[man-midengyuan]{满幂等元}[full idempotent]。例如 \(M_n(k)\) 中取 \(e=E_{11}\)，角代数就是 \(k\)；上三角矩阵代数和 \(k[t]/(t^m)\) 则已经是基本代数。基本性并不要求根为零。

\begin{proposition}\label{b2:prop:full-corner-radical-powers}
设 \(k\) 为域、\(A\) 为有限维结合含幺 \(k\) 代数，\(J=J(A)\)，\(e\in A\) 为满足 \(AeA=A\) 的幂等元，\(B=eAe\)。对每个整数 \(r\ge0\)，有
\[
J(B)^r=eJ^re.
\]
\end{proposition}
\begin{proof}
由\cref{b2:prop:corner-radical}，\(J(B)=eJe\)，而 \(r=0\) 时两边都等于 \(eAe\)。设 \(r\ge1\) 且已知 \((eJe)^r=eJ^re\)。显然 \((eJ^re)(eJe)\subseteq eJ^{r+1}e\)。反过来，由 \(AeA=A\) 可写 \(1=\sum_t a_teb_t\)。对 \(x\in J^r\)、\(y\in J\)，有
\[
exye=\sum_t(exa_te)(eb_tye)\in(eJ^re)(eJe).
\]
所以 \(eJ^{r+1}e\subseteq(eJ^re)(eJe)\)，归纳即得结论。
\end{proof}

对分裂基本代数，半单商 \(k^s\) 的坐标幂等元给出顶点，\(J/J^2\) 则记录顶点间的一次连接。一般基本代数的半单商仍可能含有非平凡除环，此时仅以普通 \(k\)-箭头记录维数会丢失顶点除环和箭头的双模结构。

\begin{definition}\label{b2:def:basic-algebra-quiver}
设 \(k\) 为域、\(B\) 为有限维分裂基本 \(k\) 代数，即 \(B/J(B)\cong k^s\)。取完整正交本原幂等元组 \(e_1,\ldots,e_s\)，并记 \(J=J(B)\)。以 \(1,\ldots,s\) 为顶点，从 \(i\) 到 \(j\) 画 \(\dim_k e_j(J/J^2)e_i\) 条箭头，得到的允许重边与圈的有向图称为 \(B\) 的\term[jiantu]{箭图}[quiver]。箭头方向配合左模约定。
\end{definition}

这张图的同构类型与所选幂等元组无关。若另选完整正交本原幂等元组 \(f_1,\ldots,f_s\)，由\cref{b2:prop:corner-radical}可知两组元素模 \(J\) 后都是 \(k^s\) 的坐标幂等元；重新编号后有 \(f_i-e_i\in J\)。对 \(x\in J\)，于是 \(f_jxf_i-e_jxe_i\in J^2\)，两组选择在 \(J/J^2\) 中给出相同的分量，箭头数相同。

\begin{proposition}\label{b2:prop:basic-algebra-arrow-generators}
设 \(k\) 为域、\(B\) 为有限维基本 \(k\) 代数，\(J=J(B)\)，并假设 \(B/J\cong k^s\)。取完整正交本原幂等元组 \(e_1,\ldots,e_s\)。对每个 \(e_j(J/J^2)e_i\) 选一组 \(k\) 基，并为每个基向量选取提升 \(x\in e_jJe_i\)，则这些提升连同 \(e_1,\ldots,e_s\) 生成代数 \(B\)。对每个左 \(B\) 模 \(M\)，有向量空间分解 \(M=\bigoplus_i e_iM\)，且提升 \(x\in e_jJe_i\) 的作用给出 \(e_iM\to e_jM\) 的线性映射。
\end{proposition}
\begin{proof}
令 \(L\) 为全部箭头提升的线性生成空间，\(C\) 为它们与全部 \(e_i\) 生成的子代数。由于提升的商类构成 \(J/J^2\) 的基，有 \(J=L+J^2\)。用 \(L^r\) 表示 \(r\) 个 \(L\) 中元素乘积的线性生成空间；展开 \(r\) 个这样的和的乘积，得
\[
J^r\subseteq L^r+J^{r+1}\subseteq C+J^{r+1}.
\]
又因 \(B/J=k^s\) 由 \(e_i\) 的商类生成，有 \(B=C+J\)。逐次把余项从 \(J\) 推到 \(J^2,J^3,\ldots\)，最终用 \(J^N=0\) 得 \(B=C\)。

对模块，\(\sum_i e_i=1\) 给出 \(m=\sum_i e_im\)，正交性使这项表达唯一，所以是向量空间直和。若 \(m=e_im\)，则 \(xm=e_jxe_im\in e_jM\)，得到所述线性映射。一般的 \(e_iM\) 未必是 \(B\) 子模；保存所有箭头的作用才能恢复各坐标之间的联系。
\end{proof}

例如 \(T_3(k)\) 的 \(J/J^2\) 有基 \(\overline E_{12},\overline E_{23}\)，给出箭头 \(2\to1\)、\(3\to2\)。元素 \(E_{13}=E_{12}E_{23}\) 来自两条箭头的复合，不再需要第三条独立箭头。一般代数中，不同路径的乘积还可能满足关系，所以箭图本身并不总能确定代数。第五卷第9.1节将引入路径代数，第9.5节进一步讨论基本代数与关系；本卷\chapref{b2:ch:16}先说明满幂等元的角代数如何与原代数的模范畴联系。


\begin{chaptersummary}
有限维性使根的幂构成的降链稳定，Nakayama 引理再使稳定项为零；半单商则把各个根滤过层变成半单模。对有限维模，\(\operatorname{rad}M=JM\)，而基座由被 \(J\) 零化的元素组成。两列的终止步数相同，但各层的具体位置和重数可能不同。

\ProseParagraphBreak
投射覆盖需要投射来源及多余核。幂等元提升把半单商中的列模转化为不可分解投射模；从最大半单商提升映射，再用 Nakayama 引理恢复满射性，便得到任意有限维模的覆盖。两个覆盖在目标上同构，不过具体同构未必唯一。简单模决定不可分解投射模的类型，简单自同态除环决定相应的重数计数。

\ProseParagraphBreak
选取每种不可分解投射类型的一份，给出基本满角代数。它的半单商去除了重复的矩阵大小；因所选幂等元满足 \(AeA=A\)，角代数的根幂还满足 \(J(eAe)^r=eJ(A)^re\)。分裂基本代数的 \(J/J^2\) 可以用箭头记录，较高次乘积则对应路径及关系；这些构造将分别在下一章的 Morita 理论和第五卷的箭图表示中继续展开。
\end{chaptersummary}

\BookExercises

\begin{exercise}\label{b2:ex:15-block-triangular-radical}
令 \(A\) 为所有分块矩阵 \(\begin{psmallmatrix}B&X\\0&C\end{psmallmatrix}\) 组成的代数，其中 \(B\in M_r(k)\)、\(C\in M_s(k)\)、\(X\in M_{r\times s}(k)\)，且 \(r,s\ge1\)。求 \(J(A)\)、半单商和根的幂零指数，并说明根以外仍可有非零幂零元素。
\end{exercise}
\begin{solution}
仅右上块可能非零的矩阵组成理想 \(I\)，任意两个这样的矩阵相乘为零，所以 \(I^2=0\)。商为 \(M_r(k)\times M_s(k)\)，是半单代数；由\cref{b2:prop:nilpotent-ideal-in-radical}及\cref{b2:prop:radical-quotient}，\(J(A)=I\)。因为 \(r,s\ge1\)，\(I\ne0\)，幂零指数为二。若 \(r\ge2\)，在左上块放 \(E_{12}\)、其余块为零，便得到不属于 \(J(A)\) 的非零平方零元素；\(s\ge2\) 时同样可用右下块。若 \(r=s=1\)，所有幂零元素恰在严格上三角理想内。
\end{solution}

\begin{exercise}\label{b2:ex:15-loewy-length-bounds}
设 \(M\ne0\) 为有限维左 \(A\) 模。证明其 Loewy 长度不超过 \(\dim_kM\)。再证明有限直和的 Loewy 长度等于各分量 Loewy 长度的最大值（全零时为零），而子模与商模的 Loewy 长度都不超过原模。
\end{exercise}
\begin{solution}
若 \(J^rM=J^{r+1}M\)，将 Nakayama 引理用于有限生成模 \(J^rM\)，得到它为零。因此到达零以前，每个相邻根层都使维数至少下降一，步数至多 \(\dim_kM\)。

对直和，\(J^r(\bigoplus_iM_i)=\bigoplus_iJ^rM_i\)，所以第 \(r\) 步为零恰好等于每个分量在该步为零，得到最大值公式。若 \(N\le M\)，则 \(J^rN\subseteq J^rM\)；商中有 \(J^r(M/N)=(J^rM+N)/N\)。原模的某次幂为零时，这两处也为零，证明所需界。
\end{solution}

\begin{exercise}\label{b2:ex:15-associated-radical-algebra}
令 \(A=T_3(k)\)、\(J=J(A)\)。计算根滤过的伴随分次代数 \(\operatorname{gr}_J(A)\)，并说明它同构于给 \(E_{ij}\) 赋次数 \(j-i\) 的原上三角矩阵代数。
\end{exercise}
\begin{solution}
对上三角矩阵，\(J^r\) 由满足 \(j-i\ge r\) 的 \(E_{ij}\) 生成；所以第 \(r\) 个分次片以 \(j-i=r\) 的矩阵单位的商类为基。把 \(E_{ij}\) 送到相应商类，是分次向量空间同构。非零乘积 \(E_{ij}E_{j\ell}=E_{i\ell}\) 的次数正好相加，其他矩阵单位乘积为零，故这也是代数同构。
\end{solution}

\begin{exercise}\label{b2:ex:15-radical-socle-subquotient}
设 \(k\) 为域。令 \(A=T_3(k)\)、\(P_i=AE_{ii}\)、\(M=P_2\oplus P_3\)，并取子模
\[
N=kE_{12}\oplus kE_{13}=\operatorname{Soc}M.
\]
计算 \(\operatorname{rad}N\)、\(\operatorname{rad}(M/N)\) 与 \(\operatorname{Soc}(M/N)\)。比较 \(N\cap\operatorname{rad}M\) 与 \(\operatorname{rad}N\)，以及商映射 \(M\to M/N\) 所诱导的 \(\operatorname{Soc}M\to\operatorname{Soc}(M/N)\) 是否满射。
\end{exercise}
\begin{solution}
记 \(J=J(A)\)。由\eqref{b2:eq:triangular-projective-layers}取 \(r=1\) 的计算，
\[
JM=kE_{12}\oplus(kE_{13}+kE_{23}),\qquad JN=0.
\]
所以 \(\operatorname{rad}N=0\)，而 \(N\cap\operatorname{rad}M=N\ne0\)。由\cref{b2:prop:radical-socle-subquotients}，商模的根为
\[
\operatorname{rad}(M/N)=(JM+N)/N=k(E_{23}+N).
\]
这里第一分量 \(P_2/kE_{12}\) 是简单模 \(S_2\)；第二分量 \(P_3/kE_{13}\) 的根为 \(k(E_{23}+N)\)，顶层为 \(S_3\)。具体地，商模中 \(E_{22}+N\) 和 \(E_{23}+N\) 都被 \(J\) 零化，而 \(E_{23}(E_{33}+N)=E_{23}+N\ne0\)。因此
\[
\operatorname{Soc}(M/N)=k(E_{22}+N)\oplus k(E_{23}+N)\cong S_2^{\oplus2}.
\]
原基座就是 \(N\)，在商中像为零，故基座上的诱导映射不满射。取商会产生原模中并非简单子模的新的基座部分；对子模有交公式，并不意味着对商也能直接取原基座的像。
\end{solution}

\begin{exercise}\label{b2:ex:15-dual-numbers-loewy}
令 \(A=k[\varepsilon]/(\varepsilon^2)\)，\(M=A^{\oplus r}\oplus k^{\oplus s}\)，其中 \(r,s\ge0\)，\(\varepsilon\) 在 \(k\) 上作用为零。计算模的根、基座、最大半单商及 Loewy 长度，并说明在 \(r>0\) 时根列不可能作为简单子模的直和分解实现。
\end{exercise}
\begin{solution}
\(J=(\varepsilon)\)，故 \(JM=(k\varepsilon)^{\oplus r}\oplus0\)，维数为 \(r\)。被 \(J\) 零化的部分是 \((k\varepsilon)^{\oplus r}\oplus k^{\oplus s}\)，故基座维数为 \(r+s\)；商 \(M/JM\cong k^{\oplus(r+s)}\)。若 \(r>0\)，Loewy 长度为二；若 \(r=0,s>0\)，为一；全零时为零。\(r>0\) 时，\(\varepsilon\) 在 \(M\) 上非零，而在任何简单模直和上均为零，所以不能把各层拼成 \(M\) 的简单直和分解。
\end{solution}

\begin{exercise}\label{b2:ex:15-change-of-scalars-truncated}
设 \(k\) 为域、\(m\ge2\)、\(A=k[t]/(t^m)\)。分别把 \(A\) 看作 \(k[t]\) 模和左 \(A\) 模。比较两种模的有限展示性、投射性与平坦性；求它们的子模、合成因子及模的根。哪些结论依赖标量环，哪些只由这里共同的子模格决定？
\end{exercise}
\begin{solution}
作为 \(k[t]\) 模，\(A\) 有展示
\[
0\longrightarrow k[t]\xrightarrow{\,t^m\,}k[t]\longrightarrow A\longrightarrow0,
\]
所以有限展示。它不投射：整环 \(k[t]\) 上的投射模是自由模的直和因子，因而无挠，但 \(\bar1\ne0\) 被 \(t^m\) 零化。它也不平坦：把单射 \(k[t]\xrightarrow{t^m}k[t]\) 与 \(A\) 张量后，得到 \(A\xrightarrow{0}A\)，不再单射。作为左 \(A\) 模，它是秩一自由模，因此有限展示、投射且平坦。

\(k[t]\) 对 \(A\) 的作用经过商映射 \(k[t]\to A\)，所以两种模的子模完全相同。它们恰为 \((\bar t^r)\)、\(0\le r\le m\)，形成一条严格下降链；每个相邻商都同构于 \(k\)，在 \(k[t]\) 模中即 \(k[t]/(t)\)，在 \(A\) 模中即 \(A/(\bar t)\)。两种模都只有一个极大子模 \((\bar t)\)，故根相同。标量环改变了投射性和平坦性，却没有改变此例的子模格、合成因子层数与根。
\end{solution}

\begin{exercise}\label{b2:ex:15-triangular-column-submodules}
求 \(T_3(k)\) 的自然列向量模 \(V=k^3\) 的全部子模，并判断它是否不可分解。
\end{exercise}
\begin{solution}
若子模 \(W\) 含向量 \(v\)，对角矩阵单位 \(E_{ii}\) 使其每个坐标向量 \(v_ie_i\) 也属于 \(W\)，因此 \(W\) 是若干坐标轴的直和。若 \(e_j\in W\)，上三角矩阵单位 \(E_{ij}\)、\(i<j\) 又使 \(e_i\in W\)。所以只有
\[
0,\quad ke_1,\quad ke_1+ke_2,\quad V
\]
这四个子模。任意两个非零子模都包含 \(ke_1\)，不可能交为零，因此 \(V\) 不可分解。它的三层依次为 \(S_3,S_2,S_1\)，与\eqref{b2:eq:triangular-projective-layers}相符。
\end{solution}

\begin{exercise}\label{b2:ex:15-triangular-projective-hom}
仍令 \(A=T_3(k)\)、\(P_i=AE_{ii}\)。计算全部 \(\operatorname{Hom}_A(P_i,P_j)\)，给出非零同态的公式，并解释为什么 \(\operatorname{End}_A(P_3)=k\) 不表示 \(P_3\) 简单。
\end{exercise}
\begin{solution}
同态由 \(E_{ii}\) 的像 \(x\) 决定，且必须满足 \(x=E_{ii}x\in E_{ii}AE_{jj}\)。反过来，这样的 \(x\) 给出 \(aE_{ii}\mapsto ax\)；因为 \(x=E_{ii}x\)，该公式与表达方式无关并且左线性。\(E_{ii}AE_{jj}\) 在 \(i\le j\) 时由 \(E_{ij}\) 生成，否则为零。因此 Hom 空间的维数表为
\[
\begin{pmatrix}1&1&1\\0&1&1\\0&0&1\end{pmatrix},
\]
第 \((i,j)\) 项的非零同态都是右乘 \(\lambda E_{ij}\) 的映射。特别地各 \(\operatorname{End}_A(P_i)=k\)。然而 \(P_3\) 有真子模 \(kE_{13}\)，所以不简单；Schur 引理的逆命题在这里不成立。
\end{solution}

\begin{exercise}\label{b2:ex:15-idempotent-nonnilpotent-counterexample}
设整数 \(n\ge2\) 有分解 \(n=\prod_{i=1}^r p_i^{a_i}\)，其中 \(p_i\) 为不同素数、\(a_i\ge1\)。分类 \(\mathbb Z/n\mathbb Z\) 的全部幂等元，求其个数，并判断哪些能提升为 \(\mathbb Z\) 的幂等元。以 \(n=12\) 具体核验，再说明“全部商幂等元可提升”是否迫使理想 \(n\mathbb Z\) 幂零。
\end{exercise}
\begin{solution}
由第一卷定理10.4.3的中国剩余定理，
\[
\mathbb Z/n\mathbb Z\cong\prod_{i=1}^r\mathbb Z/p_i^{a_i}\mathbb Z.
\]
在模 \(p^a\) 的商中，幂等条件是 \(p^a\mid e(e-1)\)。相邻整数互素，故 \(p\) 不会同时整除 \(e\) 与 \(e-1\)，整个素数幂必须整除其中一个。因此每个坐标的幂等元只有零或一。各坐标可独立选取，全部幂等元由 \(\{0,1\}^r\) 参数化，共有 \(2^r\) 个。

整数环只有零、一两个幂等元，所以可提升的商幂等元恰为各坐标全零或全一者。于是全部可提升当且仅当 \(r=1\)，也就是 \(n\) 为素数幂。模十二时，四个幂等元为 \(0,1,4,9\)：后两者分别满足 \(4^2-4=12\)、\(9^2-9=72\)，并分别在模四、模三的坐标上取 \((0,1)\)、\((1,0)\)，不能提升。

另一方面，无论 \(n\ge2\) 取何值，都有 \((n\mathbb Z)^m=n^m\mathbb Z\ne0\) 对每个正整数 \(m\) 成立。当 \(n\) 为素数幂时全部幂等元可提升，但这个理想仍不幂零。因此幂零性是正文提升定理的充分假设，不能由某个具体商中的全部幂等元可提升反推出它。
\end{solution}

\begin{exercise}\label{b2:ex:15-square-zero-matrix-lift}
令 \(R=M_2\bigl(k[\varepsilon]/(\varepsilon^2)\bigr)\)、\(I=\varepsilon M_2(k)\)。对
\[
a=E_{11}+\varepsilon\begin{pmatrix}u&v\\w&z\end{pmatrix},
\]
用平方零修正公式求出幂等提升 \(e\)。再将它写成 \((1+\varepsilon X)E_{11}(1+\varepsilon X)^{-1}\)。
\end{exercise}
\begin{solution}
相乘得 \(d=a^2-a=\varepsilon\operatorname{diag}(u,-z)\)。使用\cref{b2:thm:idempotent-lifting}中平方零情形的修正 \(e=a+(1-2a)d\)。因 \(\varepsilon^2=0\)，有 \((a-E_{11})d=0\)，故 \((1-2a)d=(1-2E_{11})d=\varepsilon\operatorname{diag}(-u,-z)\)，得到
\[
e=E_{11}+\varepsilon\begin{pmatrix}0&v\\w&0\end{pmatrix}.
\]
取 \(X=\begin{psmallmatrix}0&-v\\w&0\end{psmallmatrix}\)，则 \(XE_{11}-E_{11}X=\begin{psmallmatrix}0&v\\w&0\end{psmallmatrix}\)。又 \((1+\varepsilon X)^{-1}=1-\varepsilon X\)，展开共轭即得 \(e\)。所有计算对任意特征成立。
\end{solution}

\begin{exercise}\label{b2:ex:15-compatible-idempotent-lifts}
设 \(k\) 为域，\(R=M_2(k[\varepsilon]/(\varepsilon^2))\)、\(I=\varepsilon M_2(k)\)。求 \(E_{11},E_{22}\in R/I\) 的全部幂等元提升，并确定哪些有序提升对 \((e,f)\) 正交且满足 \(e+f=1\)。若 \(k=\mathbb F_q\)，求这样的完整正交提升组有多少个。
\end{exercise}
\begin{solution}
写 \(e=E_{11}+\varepsilon\begin{psmallmatrix}u&v\\w&z\end{psmallmatrix}\)。直接计算得
\[
e^2-e=\varepsilon\begin{pmatrix}u&0\\0&-z\end{pmatrix}.
\]
所以所有幂等提升恰为 \(e=E_{11}+\varepsilon(vE_{12}+wE_{21})\)。同样，\(E_{22}\) 的全部幂等提升为 \(f=E_{22}+\varepsilon(cE_{12}+dE_{21})\)。两方向的乘积分别为
\[
ef=\varepsilon(v+c)E_{12},\qquad fe=\varepsilon(w+d)E_{21}.
\]
于是正交恰好要求 \(c=-v,d=-w\)，此时 \(f=1-e\)，总和也为一。反过来，总和为一也给出这两个条件。因此完整正交组由任意 \((v,w)\in k^2\) 唯一确定；有限域情形共有 \(q^2\) 组。这里不仅检查单个幂等方程，还须同时满足两方向的零乘积条件。
\end{solution}

\begin{exercise}\label{b2:ex:15-superfluous-ambient}
证明多余子模的任意子模仍在同一母模中多余。若 \(N\ll P\) 且 \(N\subseteq Q\subseteq P\)，是否必有 \(N\ll Q\)？给出有限维代数中的反例。
\end{exercise}
\begin{solution}
若 \(N'\subseteq N\) 且 \(L+N'=P\)，则 \(L+N=P\)，由 \(N\ll P\) 得 \(L=P\)。所以 \(N'\ll P\)。

缩小母模却未必保持多余性。取 \(A=k[\varepsilon]/(\varepsilon^2)\)、\(P=A\)、\(N=Q=k\varepsilon\)。\(N=JP\)，由\cref{b2:prop:finite-superfluous-radical}，\(N\ll P\)。但在 \(Q=N\ne0\) 中，\(0+N=Q\) 而 \(0\ne Q\)，故 \(N\not\ll Q\)。
\end{solution}

\begin{exercise}\label{b2:ex:15-truncated-projective-cover}
令 \(A=k[t]/(t^n)\)、\(1\le s\le n\)。构造 \(M=A/(\bar t^s)\) 的投射覆盖，并将其核识别为一个循环商模。分别求来源、目标和非零核的 Loewy 长度。
\end{exercise}
\begin{solution}
自然满射 \(A\to M\) 的来源自由，核 \((\bar t^s)\subseteq JA\)，所以是覆盖。映射 \(A\to(\bar t^s)\)、\(a\mapsto a\bar t^s\) 的核为 \((\bar t^{n-s})\)，对截断多项式逐项比较系数即可看出；因此核模同构于 \(A/(\bar t^{n-s})\)。\(s=n\) 时此模为零。来源与目标的 Loewy 长度分别为 \(n,s\)，核非零时为 \(n-s\)。覆盖保留的是目标的最小投射来源，并不要求来源和目标的长度相同。
\end{solution}

\begin{exercise}\label{b2:ex:15-matrix-algebra-covers}
令 \(A=M_m(k)\)、\(S=k^m\)、\(M=S^{\oplus r}\)，其中 \(m\ge1,r\ge0\)。确定 \(M\) 的投射覆盖，并判断它何时是自由 \(A\) 模。
\end{exercise}
\begin{solution}
\(A\) 半单，\(S\cong AE_{11}\) 投射，所以 \(M\) 本身投射。恒等映射 \(M\to M\) 的核为零，是投射覆盖。正则模按列分解为 \(A\cong S^{\oplus m}\)，因而 \(m\mid r\) 时 \(M\cong A^{r/m}\)。反过来，若 \(M\) 自由，其有限 \(k\) 维数使自由基有限，设有 \(d\) 个，则 \(mr=m^2d\)，故 \(m\mid r\)。这包含 \(r=0\) 的情形，并给出许多不能选为自由模的投射覆盖。
\end{solution}

\begin{exercise}\label{b2:ex:15-cover-automorphisms}
令 \(A=k[t]/(t^n)\)、\(n\ge1\)，\(p:A\to k\) 为自然覆盖。求全部满足 \(pu=p\) 的左模自同构 \(u:A\to A\)。若 \(k=\mathbb F_q\)，求这样的自同构个数。
\end{exercise}
\begin{solution}
自同态由 \(u(1)=a\) 决定，公式为 \(u(x)=xa\)。条件 \(pu=p\) 等价于 \(a\) 的常数项为一，即 \(a\in1+(\bar t)\)。这些元素均可逆，所以它们都给出自同构，也没有别的可能。域有 \(q\) 个元素时，\(\bar t,\ldots,\bar t^{n-1}\) 的系数各有 \(q\) 种选择，个数为 \(q^{n-1}\)。因此覆盖的唯一同构类型不意味着覆盖上的自同构只能是恒等。
\end{solution}

\begin{exercise}\label{b2:ex:15-lifting-between-covers}
设 \(k\) 为域、\(A=k[t]/(t^4)\)，对 \(1\le i\le4\) 令 \(M_i=A/(\bar t^i)\)，自然投射覆盖记为 \(p_i:A\to M_i\)。固定 \(1\le r,s\le4\)，描述所有同态 \(h:M_s\to M_r\)，并求使 \(p_rH=hp_s\) 的全部 \(H:A\to A\)；判断这些提升何时可逆。分别考察 \(M_1\to M_2\)、\(1\mapsto\bar t\)，以及自然商映射 \(M_2\to M_1\)，说明目标间的映射不是同构时，提升可能具有怎样不同的行为。
\end{exercise}
\begin{solution}
同态 \(h\) 由 \(z=h(1)\in M_r\) 决定。关系 \(\bar t^s\cdot1=0\) 要求 \(\bar t^sz=0\)，即
\[
z\in\bar t^{\max(r-s,0)}M_r.
\]
比较截断多项式的系数即可得到这一条件。每个来源上的自同态形如 \(H(x)=xa\)，其中 \(a=H(1)\in A\)。提升条件恰为 \(a+(\bar t^r)=z\)；固定一个代表元 \(a_0\) 后，全部提升由 \(a\in a_0+(\bar t^r)\) 给出。这样的 \(H\) 可逆，当且仅当 \(a\) 的常数项非零：常数项为零时它在模根后为零，不能可逆；常数项非零时可用有限几何和求逆。因为 \(r\ge1\)，所有提升的常数项相同，所以对给定 \(h\)，全部提升同时可逆或同时不可逆。

对 \(M_1\to M_2\)、\(1\mapsto\bar t\)，此映射单射，所有提升都有 \(a\equiv\bar t\pmod{\bar t^2}\)，常数项为零，故都不可逆。对自然满射 \(M_2\to M_1\)，所有提升满足 \(a\equiv1\pmod{\bar t}\)，故都可逆。这两个目标映射均非同构；正文的同构提升定理给出充分条件，不能把它倒过来作为必要条件。
\end{solution}

\begin{exercise}\label{b2:ex:15-triangular-idempotents}
求 \(T_2(k)\) 的全部幂等元，判断哪些本原，并证明所有非平凡者分别共轭于 \(E_{11}\) 或 \(E_{22}\)。
\end{exercise}
\begin{solution}
写幂等元为 \(\begin{psmallmatrix}a&b\\0&c\end{psmallmatrix}\)，方程为 \(a^2=a,c^2=c,(a+c-1)b=0\)。域上的前两式使 \(a,c\in\{0,1\}\)。若两者同为零或同为一，必有 \(b=0\)；其余两种情形中 \(b\) 任意。因此全部幂等元为 \(0,1,E_{11}+bE_{12},E_{22}+bE_{12}\)。

后两族模根后是 \(k^2\) 中的坐标幂等元，由\cref{b2:prop:primitive-projective-cover}本原；\(1=E_{11}+E_{22}\) 不本原，零不计本原。又
\[
(1+cE_{12})E_{11}(1-cE_{12})=E_{11}-cE_{12},
\]
对 \(E_{22}\) 的对应公式为 \(E_{22}+cE_{12}\)。分别取 \(c=-b\) 或 \(c=b\)，就得到所需共轭。
\end{solution}

\begin{exercise}\label{b2:ex:15-corner-filtration}
设 \(k\) 为域、\(n\ge1\)。令 \(A=T_n(k)\)，选取 \(1\le q_1<\cdots<q_s\le n\)，并令 \(e=\sum_{a=1}^sE_{q_aq_a}\)。识别 \(B=eAe\)，对所有 \(r\ge0\) 分别计算 \(J(B)^r\) 与 \(eJ(A)^re\)，并判断 \(e\) 何时为满幂等元。再取 \(n=4\)、\((q_1,q_2,q_3)=(1,3,4)\)，比较二次和三次根幂。
\end{exercise}
\begin{solution}
角代数的基为 \(E_{q_aq_b}\)、\(a\le b\)，乘法保持矩阵单位规则，故按所选指标的次序同构于 \(T_s(k)\)。在 \(T_n(k)\) 中，\(J(A)^r\) 由距主对角线至少 \(r\) 的矩阵单位组成；一条长度为 \(r\) 的非零乘积，恰需选择 \(r\) 个严格增加的指标间隔。因此
\[
\begin{aligned}
J(B)^r&=\operatorname{span}_k\{E_{q_aq_b}\mid b-a\ge r\},\\
eJ(A)^re&=\operatorname{span}_k\{E_{q_aq_b}\mid q_b-q_a\ge r\}.
\end{aligned}
\]
对 \(r=0\)，两式均包含对角项，给出 \(B\)；空指标集按约定给零。前式在所选指标中数可连接的步数，后式允许在原代数中经过未选的中间指标。

若漏选某个 \(i\)，任何 \(aeb\) 的第 \(i\) 个对角元都为零，所以 \(AeA\) 不含 \(E_{ii}\)。故 \(e\) 满，当且仅当所有指标均被选中，即 \(e=1\)。在指定的四阶例子中，
\[
J(B)^2=kE_{14},\quad eJ(A)^2e=kE_{13}+kE_{14},\qquad
J(B)^3=0,\quad eJ(A)^3e=kE_{14}.
\]
虽然一次根取角的公式仍成立，较高根幂可以在不同次数都出现差异。
\end{solution}

\begin{exercise}\label{b2:ex:15-basic-corner-example}
令 \(A=M_2(k)\times k[t]/(t^2)\)。选择一个满幂等元 \(e\)，使 \(eAe\) 基本；求半单商，并给出正则左模中各类不可分解投射模的出现次数。
\end{exercise}
\begin{solution}
取 \(e=(E_{11},1)\)，则 \(eAe\cong k\times k[t]/(t^2)\)，根为 \(0\times(\bar t)\)，半单商为 \(k\times k\)，所以基本。又
\[
e+(E_{21},0)e(E_{12},0)=(I_2,1),
\]
说明 \(AeA\) 包含单位，故 \(e\) 满。两类不可分解投射模分别是 \(P_1=k^2\times0\) 与 \(P_2=0\times k[t]/(t^2)\)，作用分别经两个因子。正则模分解为 \(P_1^{\oplus2}\oplus P_2\)。第二个分量的根非零，说明选取基本角代数没有抹去这些非半单结构。
\end{solution}

\begin{exercise}\label{b2:ex:15-real-simple-division-ring}
把 \(A=M_2(\mathbb C)\) 看成实代数，并令 \(S=\mathbb C^2\) 为左模。求 \(\operatorname{End}_A(S)\) 及其在 \(\mathbb R\) 上的维数，给出 \(S\) 的投射覆盖和 \(A\) 的一个基本角代数。
\end{exercise}
\begin{solution}
任何 \(A\) 线性映射都与标量矩阵 \(iI_2\) 交换，故首先是复线性的。与 \(E_{11},E_{22}\) 交换使它为对角矩阵，再与 \(E_{12}\) 交换使两个对角元相同。因此 \(\operatorname{End}_A(S)\cong\mathbb C\)，实维数为二，在自身除环上的维数为一。

\(S\cong AE_{11}\)，且 \(A\) 半单，所以 \(S\) 投射，恒等映射是其覆盖。取 \(e=E_{11}\)，得到基本角代数 \(eAe\cong\mathbb C\)。它的半单商就是 \(\mathbb C\)，不能在实基域上改成 \(\mathbb R\)。
\end{solution}

\begin{exercise}\label{b2:ex:15-quiver-with-zero-composition}
设 \(k\) 为域、\(B=T_4(k)/(kE_{14})\)，矩阵单位的商类沿用原符号。求 \(J(B)\) 的各次幂和箭图，写出由 \(P_4=BE_{44}\) 给出的 \(S_4\) 的投射覆盖，并计算 \(P_4\) 的 Loewy 层。比较相邻箭头的二步复合与从第四顶点到第一顶点的三步复合，说明关系会在哪一层改变列模。
\end{exercise}
\begin{solution}
矩阵单位的乘法说明 \(kE_{14}\) 是双边理想。严格上三角理想的像为
\[
J(B)=kE_{12}+kE_{13}+kE_{23}+kE_{24}+kE_{34},\qquad
J(B)^2=kE_{13}+kE_{24},\quad J(B)^3=0.
\]
这是幂零理想，商为 \(k^4\)，所以它确为 Jacobson 根。\(J(B)/J(B)^2\) 的基是 \(E_{12},E_{23},E_{34}\) 的商类，箭图为 \(4\to3\to2\to1\)。两个相邻箭头的复合分别给出非零的 \(E_{24}\) 与 \(E_{13}\)，但
\[
E_{12}E_{23}E_{34}=E_{14}=0.
\]
因此该关系不消掉二步路径，而只使这条三步路径为零。

列模 \(P_4\) 的基为 \(E_{44},E_{34},E_{24}\)，根链为
\[
P_4\supset kE_{34}+kE_{24}\supset kE_{24}\supset0.
\]
自然满射 \(P_4\to S_4\) 的核是 \(J(B)P_4\)，故为投射覆盖；非零 Loewy 层依次是 \(S_4,S_3,S_2\)。原 \(T_4(k)\) 的第四列还有最底层 \(S_1\)，商掉 \(E_{14}\) 正是移除了这一层，箭头数和二步乘积均未改变。
\end{solution}

\begin{exercise}\label{b2:ex:15-projective-hom-composition-factors}
假设 \(k\) 代数闭。设 \(P_i\to S_i\) 为简单模的投射覆盖，\(M\) 为有限维模。对 \(M\) 的任一合成列，证明其中同构于 \(S_i\) 的因子个数等于 \(\dim_k\operatorname{Hom}_A(P_i,M)\)。
\end{exercise}
\begin{solution}
对简单模 \(S_j\)，所有 \(P_i\to S_j\) 的同态都通过 \(P_i/JP_i\cong S_i\) 分解。由 Schur 引理及\cref{b2:prop:split-simple-endomorphisms}，Hom 的 \(k\) 维数在 \(i=j\) 时为一，否则为零。

取合成列 \(0=M_0\subset M_1\subset\cdots\subset M_\ell=M\)。将 \(\operatorname{Hom}_A(P_i,-)\) 用于每个 \(0\to M_{r-1}\to M_r\to M_r/M_{r-1}\to0\)。\(P_i\) 投射，由\cref{b2:cor:projective-hom-exact}得到短正合列，所以这些 Hom 空间的维数逐项相加。每个因子贡献一或零，恰得到题中计数。这也表明该个数不依赖所选合成列。该计数公式可对照 Etingof 等人\cite[Proposition 9.2.3, p. 211]{EtingofEtAl2011RepresentationTheory}。
\end{solution}
